\documentclass[11pt,oneside]{article}
\usepackage[a4paper,margin=27mm,headheight=14pt]{geometry}
\usepackage[T1]{fontenc}
\usepackage[utf8]{inputenc}
\usepackage{lmodern}
\usepackage{microtype}
\usepackage{amsmath,amssymb,mathtools}
\usepackage{amsthm}
\usepackage{booktabs,longtable,tabularx,array}
\usepackage{enumitem}
\usepackage{xcolor}
\usepackage{listings}
\usepackage{fancyhdr}
\usepackage{hyperref}
\usepackage[nameinlink,noabbrev]{cleveref}
\usepackage{graphicx}
\usepackage{url}

\definecolor{scopeblue}{RGB}{24,72,112}
\definecolor{scopegreen}{RGB}{32,112,72}
\definecolor{scopeorange}{RGB}{166,87,0}
\definecolor{scopegray}{RGB}{245,247,249}
\definecolor{codegray}{RGB}{70,75,80}

\hypersetup{
  colorlinks=true,
  linkcolor=scopeblue,
  citecolor=scopegreen,
  urlcolor=scopeblue,
  pdftitle={SCOPE-iPRA: Scoped Observational Contracts and Proof Extraction},
  pdfauthor={OpenAI, prepared for Vladimir Reshetnikov}
}

\pagestyle{fancy}
\fancyhf{}
\fancyhead[L]{\small SCOPE-iPRA}
\fancyhead[R]{\small Scoped observational contracts}
\fancyfoot[C]{\thepage}

\setlength{\parindent}{0pt}
\setlength{\parskip}{0.55em}
\setlist{itemsep=0.2em,topsep=0.3em}
\renewcommand{\arraystretch}{1.15}

\newtheorem{theorem}{Theorem}[section]
\newtheorem{lemma}[theorem]{Lemma}
\newtheorem{proposition}[theorem]{Proposition}
\newtheorem{corollary}[theorem]{Corollary}
\theoremstyle{definition}
\newtheorem{definition}[theorem]{Definition}
\newtheorem{construction}[theorem]{Construction}
\newtheorem{example}[theorem]{Example}
\theoremstyle{remark}
\newtheorem{remark}[theorem]{Remark}
\newtheorem{warning}[theorem]{Warning}

\newcommand{\iPRA}{\mathrm{iPRA}}
\newcommand{\PRAAx}{\mathsf{PRAAx}}
\newcommand{\Prf}{\mathsf{Prf}}
\newcommand{\PTm}{\mathsf{PTm}}
\newcommand{\Fm}{\mathsf{Fm}}
\newcommand{\PR}{\mathsf{PR}}
\newcommand{\Scope}{\textsc{Scope-iPRA}}
\newcommand{\turnstilep}{\mathrel{\vdash_{\PRAAx}}}
\newcommand{\zero}{\mathbf{0}}
\newcommand{\one}{\mathbf{1}}
\newcommand{\bits}{\{\zero,\one\}}
\newcommand{\code}[1]{\nolinkurl{#1}}
\newcommand{\meta}{\mathrm{meta}}
\newcommand{\obj}{\mathrm{obj}}
\newcommand{\shift}{\mathord{\uparrow}}
\newcommand{\Nat}{\mathbb{N}}
\newcommand{\Bool}{\mathbb{B}}
\newcommand{\eps}{\varepsilon}
\newcommand{\res}{\mathsf{Res}}
\newcommand{\Out}{\mathsf{Out}}
\newcommand{\Valid}{\mathsf{Valid}}
\newcommand{\Obs}{\mathsf{Obs}}
\newcommand{\least}{\mathsf{least}}
\newcommand{\Search}{\mathsf{Search}}
\newcommand{\llbracket}{\mathopen{[\![}}
\newcommand{\rrbracket}{\mathclose{]\!]}}
\newcommand{\dotminus}{\mathbin{\dot{-}}}

\lstdefinelanguage{LeanMini}{
  morekeywords={structure,inductive,def,theorem,namespace,end,where,variable,open,import,by,apply,exact,match,with,if,then,else,forall},
  sensitive=true,
  morecomment=[l]{--},
  morecomment=[s]{/-}{-/},
  morestring=[b]"
}
\lstset{
  language=LeanMini,
  basicstyle=\ttfamily\small\color{codegray},
  keywordstyle=\bfseries\color{scopeblue},
  commentstyle=\itshape\color{scopegreen},
  backgroundcolor=\color{scopegray},
  frame=single,
  rulecolor=\color{scopegray},
  breaklines=true,
  columns=fullflexible,
  keepspaces=true,
  showstringspaces=false,
  xleftmargin=0.5em,
  xrightmargin=0.5em
}

\title{\textbf{SCOPE-iPRA}\\[0.4em]
\Large Scoped Observational Contracts and Proof Extraction\\
\large A general framework for embedding finitary mathematics into intuitionistic primitive recursive arithmetic}
\author{Prepared for Vladimir Reshetnikov}
\date{16 November 2026\\Repository revision \texttt{8dfe26af18310f58aaff695fb0ce03ff5007d096}}

\begin{document}
\maketitle
\thispagestyle{empty}
\vfill
\begin{center}
\begin{minipage}{0.88\textwidth}
\small
\textbf{Exact status.}  This document gives complete human proofs of the
framework theorems and a concrete Lean formalization route against the pinned
ProveIt APIs.  The structural Lean kernel in the accompanying archive has no
placeholders but was not elaborated in this runtime because no Lean executable
was available.  The bounded-reflection and compiler files are explicitly
blueprints, not claimed kernel-checked developments.  The standard-library
Python reference implementation passes all bundled regression tests; those
tests validate finite algorithms only, not object-theory derivations.
\end{minipage}
\end{center}
\vfill
\newpage

\begin{abstract}
The ProveIt iPRA corpus now contains a checked ordinary proof calculus, a large
library of checked arithmetic derivations, and more than one hundred research
packets that translate mathematics into finite quantitative claims.  The
persistent gap is no longer the slogan that ``infinite mathematics should be
made finite.''  It is a small but technically decisive proof-engineering gap:
bounded verification conditions must be reflected into quantifier-free iPRA
without losing substitution, noncanonical witnesses must be used without being
silently chosen, and the final construction must be an actual value of
\(\Prf\;\PRAAx\), with every unresolved modulus, height bound, separation gap,
or compactness datum visible as a theorem boundary.

This report proposes \Scope, a framework with three mutually reinforcing
ideas.  First, mathematical concepts are represented by \emph{query-indexed
observation presentations}: a requested finite observation, a certificate
relation, a primitive-recursive verifier, an observer, and object-iPRA
extraction and coherence derivations.  Second, bounded logic is compiled in a
\emph{schema-first} form.  A formula scheme over \(k\) formal parameters is
compiled to a fixed symbol of type \(\PR\,k\), rather than to a term that must
be recompiled after substitution.  Specialization therefore changes only the
arguments of one fixed symbol, and substitution naturality follows from the
checked syntax laws.  This removes the hardest open obligation in the earlier
bounded-reflection outlines.  Third, certificates carry an escape effect:
\emph{sealed}, \emph{relational}, or \emph{named}.  Sealed certificates are
removed by existential elimination; relational outputs expose only a coherent
observation relation; named outputs require either a direct primitive-recursive
producer or a proved finite search bound.

We prove, at the human level, a uniform bounded-reflection theorem with both
truth and falsity reflection, least bounded selection, scoped realization,
relational observer elimination, a scope-minimal normal-form theorem for finite
proof graphs, and a proof-producing compilation theorem whose output is an
open \(\Prf\;\PRAAx\) judgment.  We also give finite trace and cyclic-demand
lemmas, domain templates for analysis, algebra, geometry, probability, and
Diophantine arguments, and explicit barrier arguments explaining why exact
branch labels, global rates from metastability, or height cutoffs from mere
finiteness must not be synthesized without additional quantitative data.

The accompanying archive contains the full \LaTeX{} source and PDF, a
source-oriented Lean kernel and detailed blueprints, an executable reference
model, sixteen passing regression tests, examples, a theorem ledger, and a
source audit pinned to the current repository revision.

\end{abstract}

\tableofcontents
\newpage

\section{Repository audit and the remaining problem}
\label{sec:audit}

\subsection{The checked base is already strong}

The correct target is not an informal claim that a construction is primitive
recursive.  It is the ordinary embedded sequent
\[
   \Gamma \turnstilep \varphi,
\]
implemented by \code{Prf PRAAx}.  At the pinned revision the repository
already supplies the following checked components.

\begin{itemize}
\item \code{ConciseIPRA/Lean/ConciseIPRA/Deduction.lean} defines intuitionistic
  natural deduction with equality and proves substitution, weakening, shifting,
  simultaneous cut \code{Prf.trans}, and the usual equality conveniences.
\item \code{ConciseIPRA/Lean/ConciseIPRA/PR.lean} defines \(\PR\), \(\PTm\),
  \(\PRAAx\), and quantifier-free induction.
\item \code{Derivations/Lean/IPRADerivations/Core/Basic.lean} supplies open
  induction \code{ind} and arbitrary-term instantiation \code{inst}.
\item The checked arithmetic layer contains semiring laws, truncated
  subtraction, order, pairing, finite sums and products, bounded search,
  division, gcd, finite linear algebra, descent, and many other reusable
  results.
\item \code{Core/Definitions.lean} contains the crucial operation
  \code{PTm.toPR}: every open term is converted into the fixed primitive-
  recursive symbol that computes it, together with an object-theory equation
  \code{toPR\_eq}.  Its generic \code{recDef} supports arbitrary parameter
  arity.
\item The repository's bridge work now includes open-sequent transport between
  the ordinary presentation and the Foundation backend.  The maintained proof
  should therefore remain an ordinary \code{Prf PRAAx} derivation; a backend
  statement is a corollary rather than a competing source of truth.
\end{itemize}

These facts change the appropriate research strategy.  It is unnecessary to
invent another logic, another universal evaluator, or another monolithic proof
certificate format.  The missing framework should be a thin layer that turns
finite mathematical interfaces into values of the calculus that already
exists.

\subsection{What the packet corpus reveals}

The delivered packets repeatedly find the same finite mathematical core:
explicit rational arithmetic, finite enumerations, bounded traces, finite
covers, finite relations, and error budgets.  They also repeatedly encounter
the same failures.

\begin{enumerate}[label=\textbf{F\arabic*.}]
\item A Python check, a Mathlib theorem, or a host-level proposition is
  described as though it were an iPRA derivation.
\item A bounded quantifier is placed directly in an induction invariant, even
  though \(\PRAAx\) induction accepts only syntactically quantifier-free
  formulas.
\item Existential witnesses from different stages are silently assumed to be
  compatible, or an arbitrary witness is turned into a global function.
\item A noncanonical object---a root at collision, an eigenvector at repeated
  eigenvalue, a local chart, or a finite cover---is forced into an exact named
  representative although only invariant observations are needed.
\item An operation with quantitative error propagation is justified by mere
  negated equality rather than an apartness margin or denominator lower bound.
\item A qualitative existence theorem is used as an executable cutoff: finite
  becomes bounded height, compact becomes a supplied finite net, or convergence
  becomes a rate.
\item A proof packet gives an attractive graph compiler but leaves the generic
  Boolean reflection, binder substitution, or final \code{Prf} construction as
  an uncompiled outline.
\end{enumerate}

Recent FINER and scoped-certificate packets diagnose these issues accurately.
Their status ledgers also identify the remaining reusable kernel: generic open
bounded reflection, checked scope combinators, and a compiler that returns an
ordinary derivation.  \Scope is deliberately concentrated on that kernel.

\subsection{The design test}

A framework is successful only if it passes all of the following tests.

\begin{enumerate}[label=(\roman*)]
\item Its public endpoint is literally \(\Gamma\turnstilep\varphi\).
\item Every induction invariant is quantifier-free by construction.
\item Substitution by arbitrary compound terms is a theorem, not an informal
  convention.
\item Noncanonical certificates cannot escape their scope accidentally.
\item Exact names are created only from explicit producers or bounded search.
\item Missing global mathematics is retained as an exact residual formula.
\item The host compiler is trusted only because it constructs a kernel-checked
  \code{Prf}; its own semantic evaluation is not an axiom.
\item The framework can express relational and quotient-valued mathematics,
  not only total functions on codes.
\end{enumerate}

\section{The SCOPE thesis}
\label{sec:thesis}

\subsection{Finite observations, not completed universes}

The phrase ``embed mathematics into iPRA'' should not be read as a demand to
construct inside iPRA a completed set-theoretic universe of reals, manifolds,
measures, or categories.  A theorem makes only finitely many observations at a
time.  The appropriate object is therefore a query-indexed interface:
\[
  q \longmapsto \text{finite evidence accepted for query }q
       \longmapsto \text{finite observable output}.
\]
The theorem may quantify over all queries, but each execution and each
induction is finite.  Coherence between queries is stated explicitly when it
is needed to form a name, a limit, or a global invariant.

This observation separates two questions that are often conflated.

\begin{description}[leftmargin=2.9cm]
\item[Local verification.]  Given a finite certificate, can iPRA verify it and
  extract the finite conclusion used by the theorem?
\item[Global production.]  Can iPRA produce a certificate for every query, with
  the required bound, modulus, separation, or compatibility?
\end{description}

The first question is frequently routine and reusable.  The second may contain
the real theorem, an open quantitative problem, or an impossible selector.  A
sound framework must allow the verifier to be completed while leaving the
producer as a visible residual.

\subsection{Three layers}

\Scope separates the development into three layers.

\paragraph{Representation layer.}
A mathematical concept is given by query terms, a certificate formula or
reflected checker, an observer, and a public specification.  Infinite objects
are accessed through finite observations and coherence laws.

\paragraph{Logical layer.}
Proof-local certificates are introduced and eliminated by the ordinary
quantifier rules.  Bounded logic is reflected to characteristic bits.  The
output mode determines whether a certificate is sealed, exposed only through
an observation relation, or converted to a named term by justified selection.

\paragraph{Compilation layer.}
A fixed host-side plan combines local object derivations using substitution,
weakening, simultaneous cut, existential elimination, implication
introduction, and the reusable arithmetic modules.  Compilation returns a
\code{Prf PRAAx} value.  No object-level program interprets arbitrary proof
syntax.

\subsection{What is new relative to the earlier frameworks}

The framework makes four concrete refinements.

\begin{enumerate}[label=\textbf{N\arabic*.}]
\item \textbf{Schema-first reflection.}  Earlier outlines compile a
  \code{BForm n} to a term and then face a difficult theorem saying that
  compilation commutes with substitution under bounded binders.  Here a fixed
  formula scheme over \(k\) parameters is compiled once to a symbol
  \(\chi_A:\PR\,k\).  An instance is always \(\chi_A(\vec t)\).  Substitution
  changes \(\vec t\), not the compiler output.  Naturality is therefore the
  checked substitution law for application.
\item \textbf{Truth and falsity reflection.}  The compiler returns both
  \(\chi_A=1\leftrightarrow A\) and
  \(\chi_A=0\leftrightarrow\neg A\), as well as Booleanity.  The falsity half
  is essential for bounded universal search and for intuitionistic treatment
  of implication.
\item \textbf{Escape effects.}  Every certificate is statically classified as
  sealed, relational, or named.  This makes the presence or absence of a
  selector a type-level issue rather than a prose promise.
\item \textbf{Scope-minimal normal form.}  A finite acyclic proof graph is
  normalized so that a witness scope surrounds exactly its consumers, named
  selection is performed only at an actual escape, and unreachable producer
  boundaries are removed.  A global mixed-radix master certificate becomes an
  optional backend, not the default architecture.
\end{enumerate}

\subsection{Trust boundary}

The trusted mathematical content remains small:
\[
  \PRAAx + \text{intuitionistic first-order logic with equality}.
\]
A Lean compiler may contain bugs, but a bug cannot manufacture a proof unless
it also constructs a value accepted by the inductive \code{Prf} definition.
Reference programs, JSON traces, and exhaustive finite tests are useful for
engineering, yet they are deliberately outside the object-theory trust
boundary.

\section{Formal baseline and notation}
\label{sec:baseline}

\subsection{Terms, formulas, and derivations}

For each arity \(n\), \(\PTm\,n\) is the type of terms whose free variables are
indexed by \(\mathrm{Fin}(n)\).  Primitive-recursive symbols have types
\(\PR\,k\).  Formulas over \(n\) free variables form \(\Fm\,\PTm\,n\).
A context \(\Gamma\) is a finite list of such formulas, and
\[
   \Gamma\turnstilep A
\]
abbreviates \code{Prf PRAAx Γ A}.  De Bruijn variable \(0\) is the newest
variable.  Shifting a formula or context under a binder is written with a
superscript \(\shift\).

The following checked rules are the only structural operations required by the
core compiler.

\begin{align*}
  &\text{Substitution:} &&
  \Gamma\turnstilep A
   \Longrightarrow \Gamma[\sigma]\turnstilep A[\sigma],\\
  &\text{Weakening:} &&
  \Gamma\subseteq\Delta,\ \Gamma\turnstilep A
   \Longrightarrow \Delta\turnstilep A,\\
  &\text{Simultaneous cut:} &&
  \Gamma\turnstilep A,\quad
  (\forall B\in\Gamma)\ \Delta\turnstilep B
   \Longrightarrow \Delta\turnstilep A,\\
  &\text{Existential elimination:} &&
  \Gamma\turnstilep\exists x\,V(x),\quad
  V(b),\Gamma^{\shift}\turnstilep A^{\shift}
   \Longrightarrow \Gamma\turnstilep A.
\end{align*}

The last rule is the logical source of certificate scoping.  The conclusion
\(A\) lives at the outer arity, so the local eigenvariable \(b\) cannot occur
in it.

\subsection{Quantifier-free induction}

The induction axiom is available only for formulas recognized by
\code{Fm.QF}.  This restriction is not an obstacle when finite data are
represented by terms and bits.  A loop invariant should have the form
\[
  I(t,s,\vec p)=\one,
\]
possibly conjoined with equations and order relations on packed state terms.
Such a formula is quantifier-free even when the human interpretation of
\(I=1\) expresses a long bounded condition.  The reflection theorem later
connects that equation to the human-facing bounded formula outside the
induction.

\subsection{Open derivations are the reusable unit}

A lemma proved only for closed numerals is not an adequate library interface.
The preferred shape is an open derivation under arbitrary local hypotheses and
at arbitrary terms.  A closed generic derivation over formal variables can be
instantiated by \code{IPRA.inst}; a derivation under a context can be
specialized by \code{Prf.subst}.  This yields a simple discipline:

\begin{quote}
State every reusable arithmetic or verifier lemma at formal variables; use
substitution only at the point where a theorem-specific plan supplies terms.
\end{quote}

This discipline is particularly important for reflected formula schemes.  A
single characteristic symbol must remain valid when a query variable is later
replaced by a complicated trace state, a search result, or an encoded tuple.

\subsection{Metalevel specialization is legitimate}

A fixed theorem packet may ask Lean to construct a fixed \(\PR\) symbol from a
fixed syntax tree.  This is ordinary metaprogramming and introduces no
universal evaluator into iPRA.  The resulting symbol and its defining
object-theory equations are part of the generated proof.  What is unavailable
in general is one total primitive-recursive function that receives arbitrary
codes for all primitive-recursive programs and evaluates them.  \Scope never
requires such a function.

\section{Query-indexed presentations of mathematical concepts}
\label{sec:presentations}

\subsection{The basic presentation}

Fix an outer parameter arity \(k\), a context \(\Gamma\), and a query tuple
\(\vec q\).  A \emph{finitary observation presentation} contains the following
data.

\begin{enumerate}[label=(P\arabic*)]
\item A certificate formula \(V(\vec q,c)\), usually a fixed bounded formula or
  the equation that its characteristic bit equals \(1\).
\item An observer term \(O(\vec q,c)\), possibly a tuple code.
\item A public specification \(S(\vec q,o)\).
\item An extractor derivation
  \[
     V(\vec q,b),\Gamma^{\shift}
       \turnstilep S(\vec q^{\shift},O(\vec q^{\shift},b)).
  \]
\item Optionally, a producer
  \(\Gamma\turnstilep\exists c\,V(\vec q,c)\), or a named residual formula
  asserting exactly that producer.
\item Optionally, a coherence derivation comparing two valid certificates.
\end{enumerate}

The verifier and observer are fixed primitive-recursive constructions.  Their
mathematical interpretation may involve a finite cover, a polynomial
factorization certificate, a mesh, a path in a finite relation, a trace, a
local coordinate chart, or a rational error budget.

\subsection{Public output relations}

A noncanonical certificate naturally defines a relation rather than a
function:
\[
  \Out(\vec q,o)\;:\!\Longleftrightarrow\;
  \exists c\bigl(V(\vec q,c)\wedge o=O(\vec q,c)\bigr).
\]
This relation is a first-class mathematical interface.  It may be total,
single-valued modulo an equivalence, or merely small in diameter.  A downstream
consumer can be proved from \(\Out(\vec q,o)\) by opening the certificate
locally.  No global representative is needed.

\begin{definition}[Exact observer coherence]
An observer is exactly coherent when iPRA derives
\[
  V(\vec q,c)\wedge V(\vec q,d)
      \longrightarrow O(\vec q,c)=O(\vec q,d).
\]
\end{definition}

\begin{definition}[Quotient observer coherence]
Given a reflected equivalence relation \(\approx_{\vec q}\), the observer is
quotient coherent when
\[
  V(\vec q,c)\wedge V(\vec q,d)
      \longrightarrow O(\vec q,c)\approx_{\vec q}O(\vec q,d).
\]
\end{definition}

\begin{definition}[Metric observer coherence]
Given an integer or rational distance code and a tolerance \(e(\vec q)\), the
observer is metrically coherent when
\[
  V(\vec q,c)\wedge V(\vec q,d)
      \longrightarrow d(O(\vec q,c),O(\vec q,d))\le e(\vec q).
\]
\end{definition}

These three forms correspond respectively to exact invariants, quotient-valued
objects, and stable finite observations.

\subsection{Coherence across queries}

A family of correct answers at each precision is not yet a coherent name.  To
represent an infinite object one adds cross-query laws.  Typical forms are
\[
  q\preceq q'\ \wedge\ \Out(q,o)\ \wedge\ \Out(q',o')
   \longrightarrow d(r_{q,q'}(o'),o)\le e(q),
\]
where \(r_{q,q'}\) is a restriction or rounding map, or
\[
  d(o_q,o_{q+1})\le 2^{-q-2}
\]
with a summable motion budget.  These laws are quantitative producers, not
consequences of pointwise existence.

For a regular real name, for example, the public object is a primitive-
recursive function \(a(k)\) of rational codes together with a proved estimate
\(|a(k)-a(\ell)|\le 2^{-k}+2^{-\ell}\).  Merely proving
\(\forall k\,\exists a\,P(k,a)\) does not produce such a function.  \Scope
keeps the pointwise relation until a direct producer, bounded selector, or
compatibility theorem justifies a name.

\subsection{Compositional contracts}

Presentations compose relationally.  Suppose \(P\) produces a relation
\(R(q,x)\) and, under a local assumption \(R(q,x)\), presentation \(Q\)
produces \(T(q,x,y)\).  Their composite output is
\[
  \exists x\bigl(R(q,x)\wedge T(q,x,y)\bigr).
\]
A finite proof graph is therefore naturally a nest of existential scopes and
ordinary cuts.  Functional composition is recovered as a special case when
both relations have named primitive-recursive producers.

\subsection{Examples of presentations}

\begin{longtable}{@{}p{0.19\textwidth}p{0.32\textwidth}p{0.39\textwidth}@{}}
\toprule
Concept & Certificate & Preferred public observation \\
\midrule
Finite combinatorial object & code plus local incidence checks & exact code or invariant \\
Algebraic roots & factor/cluster and separation certificates & multiset, cluster, symmetric functions \\
Real function value & rational approximation plus error proof & interval or regular-name approximant \\
Compact optimization & finite net, coverage and value table & approximate minimum and near-minimizer relation \\
Probability law & finite rational law plus tail/quantization data & bounded expectations or robust quantile interval \\
Topological map & finite complex, carrier and homotopy checks & induced finite-chain observation \\
Diophantine solution set & height bound and checked enumeration & complete finite list \\
\bottomrule
\end{longtable}

\section{Observation modes and constructive barriers}
\label{sec:observations}

\subsection{Three escape modes}

Each certificate occurrence is assigned one of three public modes.

\begin{description}[leftmargin=2.9cm]
\item[Sealed.]  The public formula contains no certificate-dependent data.  A
  producer and extractor are composed by existential elimination.
\item[Relational.]  The public theorem mentions only an existential output
  relation or an observation modulo an explicit exact, quotient, or metric
  coherence law.  No representative is selected.
\item[Named.]  A concrete certificate code, representative, label, or observer
  term occurs in the public result.  A direct primitive-recursive producer or a
  proved finite selection bound is mandatory.
\end{description}

The distinction between relational and named output is substantive.  A
relation may be uniquely inhabited without providing a primitive-recursive
term naming its inhabitant.  Unique existence in intuitionistic arithmetic is
not a hidden description operator.

\subsection{Exact sign at zero}

Consider regular Cauchy names for real numbers and suppose a finite-query
algorithm claims to return the exact sign \(-1,0,+1\) for every name.  Run it
on the zero name.  Only finitely many approximation indices are inspected; let
\(m\) exceed all of them.  Construct a positive regular name and a negative
regular name that agree with zero at every inspected index and begin separating
only after \(m\).  The algorithm follows the same finite computation on all
three names and returns the same sign, which is incorrect for at least two.

This is not a defect of iPRA.  It is a continuity barrier for the proposed
interface.  Correct replacements are:

\begin{itemize}
\item return a sign set or interval containing all possible signs;
\item return an approximate value and postpone the branch;
\item require an apartness residual \(|x|\ge 2^{-r}\), after which a finite
  precision determines the sign;
\item formulate a branch-independent observer.
\end{itemize}

The same argument applies to an ordered root through a collision, an eigenvector
at a repeated eigenvalue, or a rotation axis at the identity.

\subsection{Metastability is not a rate}

A metastable statement has the finite-consumer form
\[
  \forall k\,\forall g\,\exists N\le B(k,g)\,
  \forall i,j\in[N,N+g(N)]\ d(a_i,a_j)\le 2^{-k}.
\]
It can be exactly what a finite proof needs.  It does not in general supply a
Cauchy rate \(R(k)\) controlling all \(i,j\ge R(k)\).  Any procedure depending
on a finite prefix can be fooled by a sequence that agrees on that prefix and
moves afterward.  To construct a regular name one must add a tail modulus, a
summable adjacent-motion bound, monotonicity plus a quantitative bound, or some
other global producer.

\subsection{Finiteness is not a height cutoff}

The assertion that a set of integer or rational solutions is finite does not
locate those solutions.  Even singleton sets can occur at arbitrarily large
codes or heights.  Bounded enumeration needs a theorem of the form
\[
  P(x)\longrightarrow h(x)\le H(\text{parameters}),
\]
or an equivalent complete-list producer.  \Scope can verify and consume such a
bound, but it does not manufacture one from cardinality alone.

\subsection{Mere nonzeroness is not an error budget}

From \(x\neq0\) intuitionistically one does not obtain a numeral \(r\) with
\(|x|\ge 2^{-r}\).  Yet a reciprocal approximation, normalization, pivot, or
Newton step consumes exactly such a lower bound.  The operation should be
totalized, and the theorem should carry a quantitative guard.  Reverse demand
propagation can then calculate the input precision required from the guard.

\subsection{Barrier-aware interface selection}

The framework uses the following decision rule.

\begin{enumerate}[label=(\arabic*)]
\item If the final theorem is independent of the witness, seal it.
\item If only an invariant or stable observation is needed, expose a relation
  and prove coherence.
\item If a concrete term is required, first look for a direct PR producer.
\item Otherwise require a finite bound, reflect the verifier, and select the
  least accepted code.
\item If none of these applies, retain a residual or reformulate the theorem.
\end{enumerate}

This rule prevents a large class of false ``constructive'' proofs while still
allowing every genuinely finite part of the argument to be formalized.

\section{Schema-first bounded reflection}
\label{sec:reflection}

\subsection{Why a schema rather than a term compiler}

A term-first bounded-formula compiler takes a syntax tree \(A\) with free
variables and returns a term \(\chi_A\).  If the variables are later replaced
by compound terms, one must compare two constructions:
\[
   \chi_{A[\sigma]}
   \qquad\text{and}\qquad
   \chi_A[\sigma].
\]
At bounded binders these may contain different primitive-recursive symbols,
because closure conversion has been performed before or after substitution.
They compute the same function, but proving that equality requires another
induction through the compiler and through the finite recursion.

\Scope removes this obligation by changing the interface.  A bounded formula
is a \emph{scheme} over \(k\) formal parameters.  Compilation returns one fixed
symbol
\[
   \chi_A : \PR\,k.
\]
For arbitrary terms \(\vec t\), the characteristic term is always
\(\chi_A(\vec t)\).  There is no operation that recompiles the scheme after
substitution.  Instantiation of the formula and instantiation of the bit both
use the same vector \(\vec t\).

\subsection{Reflected atomic schemes}

\begin{definition}[Atomic schema]
A reflected atomic scheme of arity \(k\) consists of
\begin{enumerate}[label=(\alph*)]
\item a formula \(A\in\Fm\,\PTm\,k\);
\item a symbol \(b_A\in\PR\,k\);
\item a closed derivation of Booleanity
\[
  \turnstilep b_A(\vec x)=\zero\ \vee\ b_A(\vec x)=\one;
\]
\item a closed derivation
\[
  \turnstilep\bigl(b_A(\vec x)=\one\bigr)\leftrightarrow A(\vec x).
\]
\end{enumerate}
\end{definition}

Equality and order atoms are obtained from the checked arithmetic library.  For
example, \code{eqb(s,t)} is \(1\) when \(s=t\) and \(0\) otherwise; the strict
order formula is an equality involving truncated subtraction and can therefore
be reflected by the same equality bit.  A theorem packet may also introduce a
new atom by supplying its own fixed PR verifier and object derivations.

\begin{lemma}[Falsity reflection from Booleanity]
\label{lem:atomic-zero}
For every reflected atom,
\[
  \turnstilep
  \bigl(b_A(\vec x)=\zero\bigr)\leftrightarrow\neg A(\vec x).
\]
\end{lemma}

\begin{proof}
If \(b_A=0\) and \(A\), truth reflection gives \(b_A=1\), contradicting
\(0\neq1\).  Conversely assume \(\neg A\).  Booleanity gives
\(b_A=0\vee b_A=1\).  The first disjunct is the goal; the second yields \(A\)
by truth reflection and hence absurdity.  This is an intuitionistic proof by an
object-level finite disjunction, not an appeal to unrestricted excluded middle.
\end{proof}

\subsection{Bounded formula schemes}

For each \(k\), let \(\mathsf{BCode}(k)\) be generated by
\[
\begin{array}{c}
  \mathsf{atom}(a),\quad A\wedge B,\quad A\vee B,\quad A\to B,
  \quad\neg A,\\[0.3em]
  \exists i<t\;A(i,\vec x),
  \qquad
  \forall i<t\;A(i,\vec x).
\end{array}
\]
In a bounded body the newest variable \(0\) is \(i\), while the old parameters
are shifted.  The human-facing denotation is the ordinary first-order formula
\[
\begin{aligned}
 \llbracket\exists i<t\ A\rrbracket
   &= \exists i\,(i<t\wedge\llbracket A\rrbracket),\\
 \llbracket\forall i<t\ A\rrbracket
   &= \forall i\,(i<t\to\llbracket A\rrbracket).
\end{aligned}
\]
No claim is made that these formulas are quantifier-free.  Only their
characteristic-bit equations are used as induction invariants.

\subsection{Arity-generic bounded search}

The checked repository search accepts one bound and one packed parameter.  The
underlying proof is already arity-independent: \code{recDef} permits any
number of parameters.  We use the following direct generalization.

\begin{theorem}[Parameterized least-zero search]
\label{thm:searchN}
Let \(p(i,\vec x)\) be a fixed primitive-recursive term.  There is a fixed
symbol \(\mu_p(N,\vec x)\) such that iPRA proves
\begin{align*}
 &\mu_p(N,\vec x)\le N,\tag{S1}\\
 &\mu_p(N,\vec x)<N
      \longrightarrow p(\mu_p(N,\vec x),\vec x)=0,\tag{S2}\\
 &i<\mu_p(N,\vec x)
      \longrightarrow p(i,\vec x)\neq0.\tag{S3}
\end{align*}
Moreover, if \(i<N\) and \(p(i,\vec x)=0\), then
\[
  \mu_p(N,\vec x)<N,\qquad
  p(\mu_p(N,\vec x),\vec x)=0,
  \qquad
  \mu_p(N,\vec x)\le i.
\]
\end{theorem}

\begin{proof}
Use the same primitive recursion as \code{Core/Search.lean}: at stage \(N+1\),
retain the old candidate unless it equals \(N\) and \(p(N,\vec x)=0\).  The
joint invariant is exactly (S1)--(S3), written with the current stage, a test
index, and the parameters as free variables.  It is a conjunction of equality,
order, implication, and negated equality, hence syntactically quantifier-free.
Quantifier-free induction on \(N\) proves the invariant.  The found theorem is
then a finite order argument: if the least candidate were after a supplied
zero, (S3) would contradict that zero.  This is precisely the existing proof
with the tuple \(\vec x\) left as separate parameters.  In Lean,
\code{PTm.toPR} closes the step term and \code{recDef} supplies the arbitrary
arity.
\end{proof}

\subsection{Compilation}

Assume fixed primitive-recursive bit operations
\(\mathsf{notb},\mathsf{andb},\mathsf{orb},\mathsf{impb}\), each with object
proofs on Boolean inputs.  Define \(\chi_A\) recursively.

\begin{align*}
 \chi_{\neg A} &= \mathsf{notb}(\chi_A),\\
 \chi_{A\wedge B} &= \mathsf{andb}(\chi_A,\chi_B),\\
 \chi_{A\vee B} &= \mathsf{orb}(\chi_A,\chi_B),\\
 \chi_{A\to B} &= \mathsf{orb}(\mathsf{notb}(\chi_A),\chi_B).
\end{align*}
For a bounded existential, put
\[
  p_A(i,\vec x)=\mathsf{notb}(\chi_A(i,\vec x)),
  \qquad m=\mu_{p_A}(t(\vec x),\vec x),
\]
and let \(\chi_{\exists i<t A}\) be the characteristic bit of \(m<t\).
For a bounded universal, search directly for a zero of \(\chi_A\):
\[
  r=\mu_{\chi_A}(t(\vec x),\vec x),
  \qquad
  \chi_{\forall i<t A}=\mathsf{eqb}(r,t(\vec x)).
\]
All of these are fixed symbols of arity \(k\), built by composition from the
body symbol, the bound's \code{PTm.toPR} symbol, and parameterized search.

\begin{theorem}[Uniform bounded reflection]
\label{thm:reflection}
For every fixed \(A\in\mathsf{BCode}(k)\), the compiler effectively constructs
\(\chi_A\in\PR\,k\) and closed iPRA derivations
\begin{align*}
 &\turnstilep \chi_A(\vec x)=0\ \vee\ \chi_A(\vec x)=1,
       \tag{R0}\\
 &\turnstilep
   \bigl(\chi_A(\vec x)=1\bigr)
       \leftrightarrow \llbracket A\rrbracket(\vec x),
       \tag{R1}\\
 &\turnstilep
   \bigl(\chi_A(\vec x)=0\bigr)
       \leftrightarrow \neg\llbracket A\rrbracket(\vec x).
       \tag{R2}
\end{align*}
The construction is primitive recursive at the metalevel and introduces no
object-language syntax evaluator.
\end{theorem}

The proof occupies the next two sections.  The important structural fact is
that the output of compilation is a \(\PR\) symbol.  Once (R0)--(R2) are
proved at formal variables, the checked \code{inst} theorem specializes them
under arbitrary contexts and arbitrary object terms.

\section{Proof of reflection for Boolean connectives}
\label{sec:connectives}

We give the intuitionistic argument in detail because this is where an
unmarked appeal to classical logic often enters a packet.  Throughout, suppose
\(a,b\) are Boolean terms and
\[
  a=1\leftrightarrow A,
  \qquad
  b=1\leftrightarrow B,
\]
with the corresponding zero/falsity equivalences supplied by
\cref{lem:atomic-zero} or the induction hypothesis.

\subsection{Negation}

Let \(\mathsf{notb}(a)=1\dotminus a\), where \(\dotminus\) is truncated
subtraction.  On the two Boolean cases,
\[
  \mathsf{notb}(0)=1,
  \qquad
  \mathsf{notb}(1)=0.
\]
Booleanity of \(\mathsf{notb}(a)\) follows by eliminating the proof
\(a=0\vee a=1\) and rewriting in each branch.

If \(\mathsf{notb}(a)=1\), the branch \(a=1\) is arithmetically impossible, so
\(a=0\), hence \(\neg A\).  Conversely, from \(\neg A\), falsity reflection
gives \(a=0\), and therefore \(\mathsf{notb}(a)=1\).  This proves truth
reflection for \(\neg A\); the zero half follows by the general Booleanity
lemma.

\subsection{Conjunction}

Define \(\mathsf{andb}(a,b)=a\cdot b\).  Four object-level cases on the two
Boolean disjunctions show that the product is Boolean.  If the product equals
one, neither factor can be zero; eliminating their Boolean cases gives
\(a=b=1\), hence \(A\wedge B\).  Conversely \(A\wedge B\) yields
\(a=b=1\) and the product equation.

This proof can be shortened with existing lemmas about products and zero, but
the four-case proof is useful as a generic implementation pattern: it invokes
only the reflected finite disjunctions and arithmetic rewriting.

\subsection{Disjunction}

Let \(\mathsf{orb}(a,b)=\mathsf{sg}(a+b)\), where \(\mathsf{sg}(0)=0\) and
\(\mathsf{sg}(Sx)=1\).  Again the four Boolean cases show that
\(\mathsf{orb}(a,b)=1\) exactly when at least one input is one.  Therefore
\[
  \mathsf{orb}(a,b)=1\leftrightarrow A\vee B.
\]
The converse direction is constructive: eliminate the given disjunction and
rewrite the corresponding bit to one.

An alternative implementation is
\(a+b\dotminus(a\cdot b)\).  The sign form is usually shorter because the
checked arithmetic library already exposes the zero/successor behavior of
\(\mathsf{sg}\).

\subsection{Implication}

Implication is the subtle connective.  Define
\[
  \mathsf{impb}(a,b)=\mathsf{orb}(\mathsf{notb}(a),b).
\]
We prove
\[
  \mathsf{impb}(a,b)=1\leftrightarrow(A\to B)
\]
without unrestricted excluded middle.

For the forward direction, assume the implication bit is one and assume
\(A\).  Truth reflection gives \(a=1\), hence
\(\mathsf{notb}(a)=0\).  The only way the disjunction bit can be one is then
\(b=1\), so \(B\).  Discharging the assumption gives \(A\to B\).

For the reverse direction, assume \(A\to B\) and eliminate the finite
Booleanity disjunction \(a=0\vee a=1\).  If \(a=0\), then
\(\mathsf{notb}(a)=1\), so the implication bit is one independently of \(b\).
If \(a=1\), truth reflection gives \(A\), the assumed implication gives
\(B\), and truth reflection for \(B\) gives \(b=1\).  Thus the implication bit
is again one.  The case split is justified by a characteristic bit already
inside iPRA; it does not assume that arbitrary formulas are decidable.

\subsection{Structural recursion and proof objects}

The host compiler recursively returns a record containing:
\[
  (\chi_A,\ \pi_A^{\mathrm{bool}},\ \pi_A^{1},\ \pi_A^{0}).
\]
At a connective it constructs a new fixed PR symbol by composition, then
combines the child derivations using ordinary natural-deduction constructors
and the checked arithmetic equations.  The recursion is over a finite Lean
datatype.  Its result is a syntax symbol and actual \code{Prf} terms, not a
semantic theorem about \(\Nat\).

\begin{proposition}[Connective fragment]
For every bounded-code expression formed from reflected atoms by
\(\neg,\wedge,\vee,\to\), the three conclusions (R0)--(R2) of
\cref{thm:reflection} hold.
\end{proposition}

\begin{proof}
Structural induction on the code.  The atom case is the stored package plus
\cref{lem:atomic-zero}; each connective case is the construction just given.
All generated terms are fixed compositions of primitive-recursive symbols.
\end{proof}

\section{Bounded quantifiers, witnesses, and least selection}
\label{sec:quantifiers}

\subsection{Bounded existential reflection}

Let \(A(i,\vec x)\) have reflected bit \(a(i,\vec x)\).  Put
\[
  p(i,\vec x)=\mathsf{notb}(a(i,\vec x)),
  \qquad
  m=\mu_p(t(\vec x),\vec x).
\]
By reflection,
\[
   p(i,\vec x)=0\quad\leftrightarrow\quad A(i,\vec x).
\]
Let \(e\) be the characteristic bit of \(m<t\).

\begin{lemma}[Existential reflection]
\label{lem:exists-reflection}
\[
 e=1\quad\leftrightarrow\quad
 \exists i\,(i<t\wedge A(i,\vec x)).
\]
\end{lemma}

\begin{proof}
Suppose \(e=1\).  Strict-order reflection gives \(m<t\).  By search property
(S2), \(p(m,\vec x)=0\), and hence \(A(m,\vec x)\).  Introduce the existential
with witness \(m\).

Conversely, suppose \(\exists i(i<t\wedge A(i,\vec x))\).  Open this witness
locally.  Reflection turns \(A(i,\vec x)\) into \(p(i,\vec x)=0\).  The found
part of \cref{thm:searchN} gives \(m<t\).  Reflect strict order to obtain
\(e=1\).  The formula \(e=1\) contains no local witness, so existential
elimination closes the proof.
\end{proof}

Booleanity of \(e\) comes from the strict-order atom; falsity reflection follows
from \cref{lem:atomic-zero}.  No quantified formula appears in the search
induction invariant.

\subsection{Bounded universal reflection}

Search for a counterexample by applying least-zero search directly to the body
bit:
\[
  r=\mu_a(t(\vec x),\vec x),
  \qquad
  u=\mathsf{eqb}(r,t(\vec x)).
\]

\begin{lemma}[Universal reflection]
\label{lem:forall-reflection}
\[
 u=1\quad\leftrightarrow\quad
 \forall i\,(i<t\to A(i,\vec x)).
\]
\end{lemma}

\begin{proof}
Assume \(u=1\), so \(r=t\).  Let an arbitrary \(i<t\) be given.  Eliminate
Booleanity of \(a(i,\vec x)\).  If it is one, truth reflection yields \(A(i)
\).  If it is zero, the found theorem gives \(r<t\), contradicting \(r=t\).
Thus \(A(i)\), and universal introduction closes the argument.

Conversely, assume every \(i<t\) satisfies \(A(i)\).  Eliminate the decidable
order alternative \(r<t\vee t\le r\).  In the first branch, search property
(S2) gives \(a(r,\vec x)=0\), while the universal premise gives \(A(r)\) and
hence \(a(r,\vec x)=1\), a contradiction.  In the second branch combine
\(r\le t\) from (S1) with \(t\le r\) to get \(r=t\), and reflect equality to
obtain \(u=1\).
\end{proof}

The universal proof uses only decidability of natural-number order and of the
reflected body bit.  It does not assert excluded middle for arbitrary
unbounded formulas.

\subsection{Completion of the reflection theorem}

\begin{proof}[Proof of \cref{thm:reflection}]
Proceed by structural induction on \(A\).  Atoms and Boolean connectives were
handled in \cref{sec:connectives}.  The existential case is
\cref{lem:exists-reflection}; the universal case is
\cref{lem:forall-reflection}.  Their characteristic terms are fixed
compositions of arity-generic search, bound symbols, and reflected order or
equality bits.  Falsity reflection follows uniformly from Booleanity and truth
reflection.  Every proof is built in the ordinary natural-deduction calculus.
\end{proof}

\subsection{Least accepted certificate}

Reflection turns any fixed bounded certificate predicate \(V(q,c)\) into a
Boolean verifier \(v(q,c)\).  Suppose a bound \(B(q)\) and a producer
\[
  \Gamma\turnstilep
  \exists c\bigl(c<B(q)\wedge V(q,c)\bigr)
\]
are available.  Define
\[
  c^*(q)=\mu_{\mathsf{notb}(v)}(B(q),q).
\]

\begin{theorem}[Least accepted code]
\label{thm:least-code}
Under the producer above, iPRA derives
\begin{align*}
  &c^*(q)<B(q),\\
  &V(q,c^*(q)),\\
  &\forall d<c^*(q)\;\neg V(q,d).
\end{align*}
\end{theorem}

\begin{proof}
Open the producer's witness \(c\).  Truth reflection gives a zero of the
complement verifier below \(B\), so the found theorem proves
\(c^*<B\) and that the complement verifier vanishes at \(c^*\).  Reflection
gives \(V(q,c^*)\).  For \(d<c^*\), search minimality says the complement
verifier is nonzero.  It is Boolean, so it must be one, hence \(v(q,d)=0\) and
falsity reflection gives \(\neg V(q,d)\).  The conclusion no longer mentions
the producer witness \(c\), so existential elimination seals it.
\end{proof}

\subsection{Selection is an upgrade, not a default}

The theorem justifies a named public term only because both a numerical bound
and accepted-code existence are present.  If the final theorem merely needs a
consequence invariant under all accepted certificates, applying
\cref{thm:least-code} is unnecessary.  One should open the existential locally,
prove the consequence, and seal it.  This avoids artificial code bounds and
canonical orders unrelated to the mathematics.

\section{Substitution coherence and open use}
\label{sec:naturality}

\subsection{Instantiation, not recompilation}

Let \(A\in\mathsf{BCode}(k)\) and let the compiler produce
\(\chi_A\in\PR\,k\).  Given arbitrary terms
\(\sigma:\mathrm{Fin}(k)\to\PTm\,n\), define
\[
  \chi_A[\sigma] := \mathsf{app}(\chi_A,\sigma),
  \qquad
  \llbracket A\rrbracket[\sigma]
     := \llbracket A\rrbracket.\mathsf{subst}(\sigma).
\]
Because the symbol is fixed, the syntax law for term application gives
\[
 \mathsf{subst}(\tau,\mathsf{app}(\chi_A,\sigma))
   = \mathsf{app}(\chi_A, i\mapsto\mathsf{subst}(\tau,\sigma_i)).
\]
There is no second compiler output to compare.

\begin{theorem}[Open reflection]
\label{thm:open-reflection}
For every context \(\Gamma\), every term substitution \(\sigma\), and every
bounded scheme \(A\), iPRA derives under \(\Gamma\)
\begin{align*}
 &\chi_A[\sigma]=0\vee\chi_A[\sigma]=1,\\
 &\bigl(\chi_A[\sigma]=1\bigr)
       \leftrightarrow\llbracket A\rrbracket[\sigma],\\
 &\bigl(\chi_A[\sigma]=0\bigr)
       \leftrightarrow\neg\llbracket A\rrbracket[\sigma].
\end{align*}
\end{theorem}

\begin{proof}
Apply the checked arbitrary-term instantiation theorem \code{IPRA.inst} to the
three closed generic derivations of \cref{thm:reflection}.  The formula
substitution and term-application simplification laws put the conclusions in
the displayed form.  Finally weaken from the empty context to \(\Gamma\).
\end{proof}

\subsection{Binders are handled once}

The denotation of a bounded body uses the standard lifted substitution under a
binder.  This law is part of the checked syntax implementation:
\[
  (\forall i\,B)[\sigma]
    = \forall i\,B[\mathsf{up}(\sigma)].
\]
The reflection compiler does not need its own binder-substitution recursion.
The body scheme has arity \(k+1\), and its characteristic symbol already takes
\((i,\vec x)\) as formal arguments.  When the outer parameters are specialized,
its argument vector becomes \((i,\sigma(\vec x))\) by ordinary application.

This is the principal technical advantage of schema-first compilation.  The
open obligation called ``arity-indexed PR abstraction and binder substitution''
in the earlier outline is split into components already present in the
repository:

\begin{itemize}
\item \code{PTm.toPR} closes each fixed term scheme to a symbol;
\item \code{recDef} permits arbitrary parameter arity;
\item \code{PTm.subst\_app} handles symbol application;
\item \code{Fm.subst\_subst}, \code{up}, and the binder lemmas handle formulas;
\item \code{IPRA.inst} transports closed derivations to arbitrary terms and
  contexts.
\end{itemize}

Only the arity-generic search proof and the structural reflection recursion
remain new object derivations.

\subsection{Naturality of a compiled proof plan}

Suppose a finite plan compiles to \(\Gamma\turnstilep G\).  For every
simultaneous substitution \(\sigma\), applying \code{Prf.subst} yields
\[
  \Gamma[\sigma]\turnstilep G[\sigma].
\]
This theorem applies to the final plan, not only to individual reflected
checks.  Therefore a theorem may be developed once at formal query parameters
and reused at trace states, least-search outputs, tuple projections, or
application-specific terms without rebuilding its proof graph.

\subsection{Universal closure}

An open theorem over parameters can be universally closed with repeated
universal introduction when the context permits it.  The framework thus does
not prove merely one theorem for each standard numeral.  It constructs a
single open object derivation, from which the ordinary universally quantified
mathematical statement follows inside iPRA.

\section{Scoped certificates and relational observers}
\label{sec:scopes}

\subsection{The fundamental sealing theorem}

Fix an outer context \(\Gamma\) over \(n\) variables, a certificate formula
\(V\) over \(n+1\) variables, and an outer goal \(G\).  The certificate is the
newest variable in \(V\).

\begin{theorem}[Scoped realization]
\label{thm:scoped-realization}
If
\[
  \Gamma\turnstilep\exists c\,V(c)
  \qquad\text{and}\qquad
  V(b),\Gamma^{\shift}\turnstilep G^{\shift},
\]
then \(\Gamma\turnstilep G\).
\end{theorem}

\begin{proof}
This is exactly the existential-elimination constructor \code{Prf.exE}.  The
arity of the second premise enforces the eigenvariable condition: \(b\) is
available inside the extractor, while \(G\) is an outer formula and cannot
contain it.
\end{proof}

The logical proof is elementary; the architectural conclusion is powerful.  A
certificate can be used many times, can feed dependent local producers, and
can support arbitrary finite reasoning.  It need not have a numerical code
bound unless it is used by an object-level variable-horizon computation or
must escape as a public name.

\begin{corollary}[Conditional boundary]
\label{cor:conditional}
From an extractor
\(V(b),\Gamma^{\shift}\turnstilep G^{\shift}\) alone, one obtains
\[
  \Gamma\turnstilep (\exists c\,V(c))\to G.
\]
\end{corollary}

\begin{proof}
Introduce the producer formula as a hypothesis, eliminate its existential, and
use the extractor after weakening by the shifted implication hypothesis.
\end{proof}

This is the correct endpoint when a height bound, cover, modulus, or analytic
constant has not yet been produced.  The verifier and all downstream
mathematics can still be formalized completely.

\subsection{Dependent scoped composition}

Suppose the second certificate depends on the first:
\begin{align*}
 &\Gamma\turnstilep\exists x\,A(x),\\
 &A(x),\Gamma^{\shift}\turnstilep\exists y\,B(x,y),\\
 &B(x,y),A(x),\Gamma^{\shift\shift}\turnstilep G^{\shift\shift}.
\end{align*}

\begin{theorem}[Dependent sealing]
\label{thm:dependent-sealing}
The three derivations compose to \(\Gamma\turnstilep G\).
\end{theorem}

\begin{proof}
Eliminate the outer existential over \(x\).  Within that scope eliminate the
existential over \(y\).  Apply the final extractor and close the scopes in
reverse order.  The shifts are exactly those already managed by the two
\code{exE} constructors.
\end{proof}

Iterating this proof realizes any finite acyclic dependency tree.  A DAG with
shared consumers is handled by opening a certificate once around the combined
continuation, or by deriving a conjunction of all required consequences inside
one scope.

\subsection{Exact relational output}

Let \(O(c)\) be an observer term and define
\[
  \Out(o)\;:\!\Longleftrightarrow\;
    \exists c\,(V(c)\wedge o=O(c)).
\]

\begin{theorem}[Existence and uniqueness of an observed relation]
\label{thm:observed-relation}
Assume a producer \(\Gamma\turnstilep\exists c\,V(c)\) and observer coherence
\[
  V(c),V(d),\Gamma^{\shift\shift}
     \turnstilep O(c)=O(d).
\]
Then iPRA proves:
\begin{enumerate}[label=(\roman*)]
\item \(\exists o\,\Out(o)\);
\item \(\Out(o)\wedge\Out(o')\to o=o'\).
\end{enumerate}
\end{theorem}

\begin{proof}
For existence, open a valid certificate and introduce the witness \(O(c)\) for
\(o\), then introduce \(c\) in the definition of \(\Out\).  The outer
conclusion contains neither local variable, so both existential eliminations
close.

For uniqueness, open the certificate hidden in each output relation.  Their
validity hypotheses permit the coherence derivation, and the two equations
identifying \(o,o'\) with the observers transfer that equality to the public
outputs.  Close both certificate scopes.
\end{proof}

The theorem proves a uniquely inhabited relation.  It does not by itself
construct a primitive-recursive term naming the unique output.  A downstream
theorem may consume the relation directly.

\subsection{Quotient and metric output}

Replace equality in the coherence premise by a reflected equivalence relation
\(\approx\).  The same proof yields
\[
  \Out(o)\wedge\Out(o')\longrightarrow o\approx o'.
\]
This is the appropriate endpoint for projective directions, multisets, roots
modulo permutation, charts modulo transition, or codes modulo observational
equality.

For a metric observer, suppose
\[
  V(c)\wedge V(d)\longrightarrow
  \operatorname{dist}(O(c),O(d))\le e.
\]
Opening two output relations gives the same diameter estimate for all public
outputs.  A finite consumer whose own modulus tolerates \(e\) can therefore use
any local representative and seal it afterward.  A canonical representative
is not required.

\subsection{Named output}

There are exactly two supported ways to export a concrete term.

\begin{enumerate}[label=(\arabic*)]
\item A direct primitive-recursive producer term \(p(q)\) together with a proof
  of \(V(q,p(q))\).
\item A finite bound and accepted-code existence, followed by the least selector
  of \cref{thm:least-code}.
\end{enumerate}

An existential eigenvariable cannot be moved into an outer term.  Any apparent
third method hides either an unproved choice principle, an omitted bound, or a
change of theorem to a relational output.

\section{Escape effects and scope-minimal normal form}
\label{sec:effects}

\subsection{Static effect information}

A proof graph node that introduces a certificate is annotated with one of the
following effects.

\begin{center}
\begin{tabularx}{0.94\textwidth}{@{}lXl@{}}
\toprule
Effect & Permitted public dependence & Required evidence \\
\midrule
\(\mathsf{sealed}\) & none & producer or residual; local extractor \\
\(\mathsf{rel}(O,\approx)\) & existential observation relation modulo \(\approx\) & producer or residual; coherence proof \\
\(\mathsf{named}(p)\) & concrete PR term \(p(q)\) & direct validity proof \\
\(\mathsf{named}(\mu,B)\) & least accepted code below \(B(q)\) & reflected verifier and bounded existence \\
\bottomrule
\end{tabularx}
\end{center}

The first two modes do not move a certificate eigenvariable to the outer
scope.  The named modes replace it with an outer term already defined without
that eigenvariable.

A host-side checker also records whether a missing producer has a named
residual.  A node with no producer and no residual is rejected; otherwise the
residual becomes an explicit context formula in the compiled theorem.

\subsection{Scope regions}

For an acyclic dependency graph, the \emph{consumer region} of a certificate
node is the subgraph consisting of all nodes whose local derivations depend on
that certificate before it is sealed or converted to a justified public
output.  Opening the witness outside this region is harmless but bloats the
proof and may force unrelated code bounds.  Opening it inside the region would
leave some consumer without the witness.

A scope-minimal plan opens the certificate immediately before the least proof
region containing all its consumers and closes it immediately after that
region's result has become certificate-independent.

\begin{theorem}[Scope-minimal normal form]
\label{thm:scope-normal-form}
Let \(G\) be a finite acyclic proof graph in which every certificate node has a
valid escape effect and every local edge is backed by an object-iPRA derivation.
Then \(G\) can be transformed effectively into a nested proof plan with the
following properties.

\begin{enumerate}[label=(\roman*)]
\item Every sealed certificate is introduced by one existential elimination
  surrounding exactly its live consumer region.
\item Every relational certificate is hidden inside its public output relation;
  no representative occurs free outside that relation's eliminator.
\item Every named certificate has been replaced by a direct producer term or a
  least bounded selector before it appears in an outer formula.
\item Unreachable nodes and their producer residuals are absent.
\item All remaining unresolved producer formulas occur in the outer context.
\end{enumerate}
\end{theorem}

\begin{proof}
Choose a reverse topological order of the graph.  Remove every node not
reachable from the final goal.  For a live certificate \(c\), collect all edges
whose derivations require its validity formula.  Replace the finite collection
of downstream conclusions by one continuation deriving their conjunction or
the final joined result under the single local assumption \(V(c)\).  This is
possible because weakening permits unrelated local assumptions and ordinary
conjunction introduction shares the witness inside one scope.

Now move the existential elimination inward across any inference whose other
premises and conclusion do not contain the certificate.  The standard
commuting conversion is valid because the eigenvariable is fresh there.  Each
move strictly decreases the number of graph edges crossed by the scope, so the
finite process terminates at the least consumer region.  Repeat in reverse
nesting order for dependent certificates.

For a relational effect, replace the certificate-dependent observer term in
the outer interface by the existential relation \(\Out\); the proofs of
\cref{thm:observed-relation} justify all public comparisons.  For a named
effect, insert the supplied producer or \cref{thm:least-code}, yielding a term
whose free variables are outer parameters only.  Finally move unresolved
producer assumptions outward by implication/context introduction.  Every
transformation is a composition of weakening, cut, conjunction rules,
existential rules, and substitution, so the resulting plan has the same
object-theory conclusion.
\end{proof}

\subsection{No-escape safety}

\begin{corollary}[No dangling certificate]
A well-effected plan cannot compile to a public formula containing a proof-local
certificate variable.
\end{corollary}

\begin{proof}
In the sealed and relational cases the certificate is bound by an existential
inside the plan.  In the named cases it is replaced by an outer term.  The
indexed type of existential elimination requires the final formula to live at
the outer arity, so no unbound local variable can occur.
\end{proof}

In Lean this property does not need a separate free-variable checker.  It is a
consequence of the indices of \code{Fm}, contexts, and \code{Prf.exE}.  The
host effect checker adds a mathematical policy: it rejects an exact public
representative when the formal scope could only support a relation.

\subsection{Master certificates as an optional backend}

If every internal witness has an explicit finite bound, the nested plan can be
packed into one mixed-radix master code.  Its checker is the conjunction of all
local checks, and bounded search selects the least accepted master code.
Conversely, a master code can be opened once and its fields used in nested
scopes.  Thus bounded scoped and packed presentations are extensionally
intertranslatable.

The scoped representation is strictly more general as an interface: it remains
meaningful when a proof-local existential witness has no useful numerical
bound.  It also avoids changing every radix when one local certificate format
changes.

\section{The proof-producing compiler}
\label{sec:compiler}

\subsection{A small indexed plan language}

The compiler does not need a universal graph interpreter.  A compact host
language suffices:

\begin{description}[leftmargin=3.0cm]
\item[Exact.] Store an existing derivation \(\Gamma\turnstilep A\).
\item[Specialize.] Apply a simultaneous term substitution.
\item[Weaken.] Enlarge the local context.
\item[Supply.] Replace every assumption of a derivation by a derivation under a
  new context, using simultaneous cut.
\item[Seal.] Combine a producer of \(\exists c\,V(c)\) with a continuation under
  \(V(c)\).
\item[Boundary.] Retain a missing producer as an assumption or implication.
\item[Reflect/search/trace.] Invoke a reusable module whose endpoint is already
  an ordinary object derivation.
\end{description}

A Lean GADT can index a plan by its arity, context, and goal:

\begin{lstlisting}
inductive Plan : {n : Nat} -> List (Fm PTm n) -> Fm PTm n -> Type
  | exact (h : Gamma |- [PRAAx] A) : Plan Gamma A
  | specialize (p : Plan Gamma A) (sigma : Fin n -> PTm m) :
      Plan (Gamma.map (fun C => C.subst sigma)) (A.subst sigma)
  | supply (p : Plan Gamma A)
      (s : forall C, C in Gamma -> Plan Delta C) : Plan Delta A
  | seal (producer : Plan Gamma (Fm.ex V))
      (consumer : Plan (V :: Gamma.map Fm.shift) A.shift) : Plan Gamma A
  | boundary (body : Plan (B :: Gamma) A) : Plan Gamma (B --> A)
\end{lstlisting}

The notation in the archive uses the repository's Unicode symbols; the display
above emphasizes the indices.

\subsection{Compilation theorem}

\begin{theorem}[Plan soundness]
\label{thm:plan-soundness}
There is a total structural function
\[
  \mathsf{compile}:\mathsf{Plan}(\Gamma,A)
     \longrightarrow (\Gamma\turnstilep A).
\]
\end{theorem}

\begin{proof}
The exact case returns its stored derivation.  Specialization is
\code{Prf.subst}; weakening is \code{Prf.weaken}; supply is
\code{Prf.trans}; seal is \code{Prf.exE}; boundary is \code{Prf.impI}.
Reflection, search, and trace constructors store the object derivations
returned by their corresponding modules.  Structural recursion terminates on
the finite host syntax.
\end{proof}

This theorem is intentionally almost tautological.  That is a virtue: the
compiler's trusted-facing kernel consists of a handful of already checked
constructors.  Sophisticated scheduling, packing, optimization, or graph
normalization may be wrong at the host level, but its output is accepted only
when it contains a valid \code{Prf} value with the declared indices.

\subsection{Graph compilation}

A graph frontend performs the following metalevel passes.

\begin{enumerate}[label=\textbf{C\arabic*.}]
\item Remove nodes unreachable from the requested public goal.
\item Check every edge's arity and local formula interface.
\item Assign and validate certificate escape effects.
\item Compute finite precision/resource demands where the graph contains
  quantitative nodes.
\item Normalize witness scopes by \cref{thm:scope-normal-form}.
\item Topologically schedule exact proof nodes and combine them by
  \code{Prf.trans}.
\item Insert reflected checks, bounded selectors, certified traces, or
  relational observer eliminators.
\item Collect reachable unresolved producer formulas in a residual context.
\item Build a \code{Plan} and call \(\mathsf{compile}\).
\end{enumerate}

No pass is an object-theory evaluator.  The graph is fixed data at the Lean
metalevel; each pass constructs fixed syntax and derivations.

\subsection{Boundary-exact endpoint}

Let \(R_1(q),\ldots,R_s(q)\) be the unresolved producer formulas reachable from
the goal after dead-node elimination.  The compiler returns
\[
   R_1(q),\ldots,R_s(q),\Gamma(q)
      \turnstilep G(q).
\]
Repeated implication introduction yields a closed conditional form if desired.
A supplied producer derivation discharges its boundary by simultaneous cut.

\begin{proposition}[Dependency exactness]
The residual list contains no producer node that is unreachable in the chosen
plan, and every reachable node without a producer contributes a residual.
\end{proposition}

\begin{proof}
The first claim follows from the initial reachability pass.  The second is the
static well-formedness rule for certificate nodes.  Compilation has no
constructor for silently assuming a producer.
\end{proof}

Dependency exactness is not logical minimality.  One residual may follow from
another by a nontrivial theorem.  Proving such an implication is ordinary
mathematics and should not be hidden in a compiler optimization.

\subsection{Proof sharing and size}

Fixed logical structure is shared by Lean definitions.  Variable-length
computation is represented by a primitive-recursive trace and one induction,
not an unrolled proof.  A connected subgraph may later be compiled to the
repository's native clause calculus for compression, provided its export back
to ordinary \code{Prf PRAAx} is checked.  These are representation
optimizations; the canonical theorem remains the ordinary sequent.

\section{Certified traces, cyclic demands, and relational paths}
\label{sec:traces}

\subsection{Deterministic trace reflection}

Let a fixed state machine have primitive-recursive initial state
\(s_0(\vec q)\), step
\(s_{t+1}=F(t,s_t,\vec q)\), and a quantifier-free invariant
\(I(t,s,\vec q)\).  The state sequence is defined by \code{recDef}.

\begin{theorem}[Trace invariant]
\label{thm:trace}
Assume object derivations
\begin{align*}
 &\Gamma\turnstilep I(0,s_0(\vec q),\vec q),\\
 &I(t,s_t,\vec q),\Gamma^{\shift}
     \turnstilep I(t+1,F(t,s_t,\vec q),\vec q).
\end{align*}
Then iPRA proves \(I(T,s_T,\vec q)\) for every term \(T\).
\end{theorem}

\begin{proof}
Use quantifier-free induction on the formula
\(I(t,s_t,\vec q)\).  The zero and successor equations of \code{recDef}
rewrite the two induction cases to the supplied derivations.  Instantiate the
open result at \(T\).
\end{proof}

A human invariant containing bounded quantifiers is first compiled to a bit
\(i(t,s,\vec q)\); induction is applied to the equation \(i=1\).  Open
reflection recovers the human statement at the terminal state.

\subsection{Query-first precision scheduling}

An analytic or geometric computation should not carry an entire infinite name
through one loop.  Fix the terminal query \(q\) first, propagate the finite
input precision and resource requirements backward, and specialize every node
to those demands.  The forward trace then manipulates only finite state.  The
final theorem may quantify over \(q\) outside the trace.

This order is logically important.  It turns ``for every precision there is a
finite computation'' into one open theorem parameterized by precision, without
placing a universal quantifier inside the induction invariant.

\subsection{Finite monotone cyclic demand closure}

A reverse demand graph may have cycles.  Let
\[
  D=\prod_{j<m}\{0,\ldots,C_j\}
\]
with componentwise order, and let \(F:D\to D\) be monotone.  Define the
inflationary step
\[
  T(d)=d\vee F(d)
\]
with componentwise maximum and clipping at capacities.  Starting from a base
demand \(d_0\), put \(d_{r+1}=T(d_r)\).

\begin{theorem}[Capped demand stabilization]
\label{thm:demand}
Let \(H=\sum_{j<m}C_j\).  Then \(d_H\) is a prefixed point of \(F\).  It is the
least prefixed point above \(d_0\).
\end{theorem}

\begin{proof}
The sequence is componentwise nondecreasing.  Define the rank
\(\rho(d)=\sum_j d_j\).  If \(d_{r+1}\neq d_r\), at least one coordinate
strictly increases, so \(\rho(d_{r+1})\ge\rho(d_r)+1\).  The rank is at most
\(H\).  Therefore there can be at most \(H-\rho(d_0)\) strict steps; by stage
\(H\), \(T(d_H)=d_H\), hence \(F(d_H)\le d_H\).

For leastness, let \(p\ge d_0\) satisfy \(F(p)\le p\).  Induct on \(r\).  If
\(d_r\le p\), monotonicity gives \(F(d_r)\le F(p)\le p\), so
\(d_{r+1}=d_r\vee F(d_r)\le p\).  Thus \(d_H\le p\).
\end{proof}

For fixed dimension, all coordinates, comparisons, maxima, and the rank are PR
terms.  The two inductions use quantifier-free vector equations and order
formulas.  An overflow bit records whether the chosen caps were adequate; its
vanishing is either proved or retained as a residual.

\subsection{Relational paths and backward viability}

Pointwise existence at successive stages does not imply that independently
chosen witnesses form a compatible path.  Let finite layers
\(X_0,\ldots,X_T\) have reflected validity predicates and reflected edge
relations \(E_t(x,y)\).  Define viability backward:
\[
\begin{aligned}
  \mathsf{viable}_T(x)&=\mathsf{accept}(x),\\
  \mathsf{viable}_t(x)&=1
   \quad\Longleftrightarrow\quad
   \exists y<B_{t+1}\,[E_t(x,y)\wedge\mathsf{viable}_{t+1}(y)=1].
\end{aligned}
\]
Each existential is a bounded reflected search.

\begin{theorem}[Least viable path]
If an initial viable node exists, bounded least search at each forward stage
constructs a whole path whose edges and terminal acceptance are verified.
Conversely, viability zero implies that no accepted path begins there.
\end{theorem}

\begin{proof}
Backward induction establishes the equivalence between the viability bit and
existence of an accepted suffix.  In the forward pass, the current node's
viability supplies an accepted successor below the finite bound; choose the
least one by \cref{thm:least-code}.  Continue for the fixed horizon.  The
terminal viability equation is exactly acceptance.  The no-path direction
opens any purported path and propagates its suffix backward, contradicting
viability zero.
\end{proof}

This theorem resolves hidden finite choice.  A greedy locally valid choice is
never assumed compatible; the whole finite interaction is verified before any
representative path is named.

\section{Reusable domain templates}
\label{sec:templates}

This section describes how common mathematical concepts fit the same small
interface.  The templates are not claims that every theorem in the domain is
easy; they separate the routine finite verifier from the theorem-specific
producer.

\subsection{Finite algebra and combinatorics}

A finite structure is encoded by a natural number or a bounded tuple.  Local
relations---incidence, adjacency, multiplication tables, color constraints,
permutation conditions---are fixed bounded formulas and hence have reflected
checkers.  The preferred workflow is:

\begin{enumerate}[label=(\arabic*)]
\item use structural pairing for a fixed record and bounded arrays for
  variable-length tables;
\item reflect the well-formedness predicate;
\item run bounded search or a certified trace;
\item state the public theorem at arbitrary input terms;
\item eliminate noncanonical auxiliary tables unless the output really needs
  them.
\end{enumerate}

Counting arguments use finite sums, injections, and the checked pigeonhole and
descent modules.  A host enumeration may suggest a certificate, but the object
proof must verify every field and derive the conclusion from the reflected
acceptance equation.

\subsection{Rational and real analysis}

Rational numbers are represented by signed integer numerators and positive
denominators, with equality and order reduced to integer polynomial
identities.  A finite real observation is an interval or rational
approximation with a stated error.

A regular-name presentation for a real \(x\) consists of a PR approximant
\(a(k)\) and an object proof
\[
  |a(k)-a(\ell)|\le 2^{-k}+2^{-\ell}.
\]
Operations are implemented query-first.  Addition propagates the desired
output precision to the two inputs.  Multiplication also requests magnitude
bounds.  Reciprocal and normalization require an apartness resource
\(|x|\ge2^{-r}\).  The reverse demand pass calculates the needed input
precision from \(k,r\).

For an existence proof that supplies only approximate witnesses, use a
relational presentation at each precision and prove cross-precision stability.
Do not infer a regular name from pointwise diameter alone.  A name requires an
explicit compatible selector or a stability theorem strong enough that the
least finite selector is coherent.

\subsection{Algebraic and multivalued data}

Polynomial roots, eigenvectors, logarithm branches, Puiseux expansions, and
local inverse branches are naturally relational near degeneracies.  The
certificate may contain:

\begin{itemize}
\item a factorization or approximate factorization;
\item rational disks or boxes containing clusters;
\item a root-count certificate on boundaries;
\item Newton-polygon data;
\item a separation gap;
\item a permutation or matching between two refinements.
\end{itemize}

At collisions, export a multiset, cluster tree, spectral projector, invariant
subspace, or symmetric polynomial observation.  Prove coherence modulo
permutation or within a metric diameter.  If a named branch is genuinely
required, add a positive gap and a deterministic ordering whose stability is
proved from that gap.  The gap is a resource, not a cosmetic side condition.

\subsection{Compactness, optimization, and topology}

Replace abstract compactness by a finite-net interface:

\begin{enumerate}[label=(\alph*)]
\item a PR list of net points at precision \(k\);
\item a reflected coverage relation;
\item a modulus transferring estimates from a point to a nearby net point;
\item bounded search for the least accepted point or least value code.
\end{enumerate}

The coverage theorem is a producer.  Once supplied, finite minimization is
routine.  If minimizers are nonunique, export the minimum value exactly and a
near-minimizer relation, or prove that all selected minimizers have close
observations.

For topology, use finite complexes, carriers, chain maps, and explicit prism
homotopies.  A local-to-global construction should descend only the observation
needed by the theorem.  A checked spanning tree can transport local data and
accumulate rational edge errors; constructing a global quotient object is
unnecessary unless it appears in the endpoint.

\subsection{Probability and statistics}

Begin with finite rational laws.  Expectation is a finite weighted sum;
variance and moments are polynomial identities after denominator clearing.
For larger spaces, a presentation must supply finite quantizations, tail
bounds, integration moduli, and any anti-concentration data used by a
discontinuous statistic.

At an atom, an exact median or quantile may be nonunique or discontinuous.
Export the interval of valid quantiles or a robust bracket.  A unique exact
quantile can be named only under a gap or anti-concentration premise that a
finite checker can consume.

\subsection{Diophantine and arithmetic geometry}

Given an explicit height bound \(H\), enumerate every code below an effective
bound \(E(H)\), test the Diophantine relation, and return the complete list.
The verifier proves:

\begin{enumerate}[label=(\roman*)]
\item every listed item is a solution;
\item no list entry is duplicated under the chosen code normalization;
\item every solution of height at most \(H\) occurs in the list.
\end{enumerate}

The remaining theorem is exactly
\(P(x)\to h(x)\le H\).  It may be a Runge estimate, a descent theorem, a Baker
bound, a supplied covering list, or an open effective producer.  Mordell--Weil
and Selmer computations likewise separate finite class arithmetic from the
existence of rational representatives on coverings.  The class verifier can
be complete while the representative producer remains a residual.

\subsection{Categories and abstract structures}

For a fixed finite presentation, categorical constructions are finite tables
with equation checkers.  Composition, finite limits and colimits, finite Kan
extensions, and finite descent can therefore be treated by bounded reflection.
The key restriction is ``fixed finite presentation.''  A universal evaluator
for arbitrary categories or arbitrary PR programs is neither needed nor
available.  Functoriality and naturality are public object derivations over the
fixed codes.

\section{Resolution of recurring difficulties}
\label{sec:difficulties}

The following matrix distinguishes four outcomes.

\begin{description}[leftmargin=2.8cm]
\item[Resolved.] A generic object-iPRA schema is supplied.
\item[Conditional.] The finite verifier is generic once a named quantitative
  producer is supplied.
\item[Reformulated.] The requested exact interface is not constructively stable;
  a relation, quotient, interval, or guard replaces it.
\item[Open producer.] The framework isolates but does not prove the deep global
  theorem.
\end{description}

\small
\begin{longtable}{@{}p{0.23\textwidth}p{0.49\textwidth}p{0.18\textwidth}@{}}
\toprule
Difficulty & SCOPE treatment & Status \\
\midrule
\endfirsthead
\toprule
Difficulty & SCOPE treatment & Status \\
\midrule
\endhead
Host computation cited as proof & Every node stores or constructs a literal
\(\Prf\;\PRAAx\) derivation; tests are outside the trust boundary & Resolved \\
Bounded quantifier in induction & Compile a fixed scheme to a Boolean PR symbol;
induct on the equation that the bit is one & Resolved \\
Binder substitution & Compile schemes to \(\PR\,k\), then instantiate only the
arguments; use checked syntax substitution & Resolved design \\
Hidden finite choice & Seal local witnesses or verify a whole path by backward
viability & Resolved boundedly \\
Arbitrary witness turned into function & Keep a relational output; name only by
direct producer or bounded least search & Resolved interface \\
Cyclic precision demands & Monotone capped closure with rank bound and overflow
residual & Resolved under caps \\
All precisions inside one loop & Fix the terminal query, propagate demands, and
run a finite specialized trace & Resolved schema \\
Exact sign/branch at collision & Export sign set, cluster, quotient observer, or
require a positive gap & Reformulated \\
Eigenvector at repeated eigenvalue & Use invariant subspace/projector or a
gap-guarded normalized vector & Reformulated \\
Division from \(x\neq0\) & Totalize and require an apartness/lower-bound resource & Conditional \\
Abstract compactness & Supply finite nets, coverage checker, and modulus & Conditional \\
Metastability promoted to rate & Keep the finite-window theorem; require a tail
modulus for a global name & Barrier \\
Pointwise approximants called a name & Add cross-precision stability or a
compatible bounded selector & Conditional \\
Qualitative finiteness used as cutoff & Require a height bound or complete-list
producer & Open producer \\
Selmer classes treated as points & Separate finite classes from representative
certificates & Open producer \\
Ineffective constants & Put thresholds and estimates in the residual list and
verify the final inequality finitely & Conditional \\
General probability space & Use finite laws or supply quantization, tails and
integration moduli & Conditional \\
Completed topological object & Use finite complexes/nets and explicit comparison
producers & Conditional \\
Code equality for equivalent objects & Use observational setoids and congruence & Resolved architecture \\
Finite gluing inflated to quotient & Use spanning-tree observational descent & Resolved finitely \\
Universal PR evaluator & Specialize fixed syntax at the metalevel & Rejected interface \\
Open theorem not portable & Maintain ordinary open \(\Prf\;\PRAAx\); transport
through the checked bridge as a corollary & Resolved architecture \\
\bottomrule
\end{longtable}
\normalsize

\subsection{Why the status distinctions matter}

A framework that labels every obstacle ``resolved'' merely moves mathematics
into unnamed assumptions.  The residual discipline prevents that.  For
example, a complete finite verifier for integral points is valuable, but it
does not prove a general effective Siegel bound.  The compiled endpoint should
say exactly:
\[
  \mathsf{HeightBound}(C,H),\ \Gamma
     \turnstilep \mathsf{CompleteIntegralPointList}(C,H,L).
\]
A later theorem may discharge \(\mathsf{HeightBound}\).  Until then the result
is a precise conditional theorem rather than an overclaim.

Likewise, reformulation is not failure.  At a multiple root the unordered
cluster is often the invariant actually used by downstream mathematics.  A
proof that works uniformly on the cluster can be stronger and more stable than
a discontinuous branch label.

\subsection{A packet audit checklist}

Before a theorem packet is described as an iPRA result, check:

\begin{enumerate}[label=$\square$]
\item The endpoint is a literal ordinary derivation.
\item Every loop invariant has an explicit \code{QF} proof.
\item Every bounded formula has truth and falsity reflection.
\item Every formula and verifier specializes at arbitrary terms.
\item Every certificate has an escape mode.
\item Every named certificate has a producer or finite bound.
\item Every multi-stage finite choice has a compatibility/path proof.
\item Every noncanonical observer has exact, quotient, or metric coherence.
\item Every partial operation has the quantitative guard it consumes.
\item Every compactness claim supplies a finite net and coverage proof.
\item Every claimed list completeness has a height/cutoff producer.
\item Semantic tests are labelled separately from object proof status.
\item The repository revision and remaining residuals are recorded.
\end{enumerate}

\section{Worked examples}
\label{sec:examples}

The examples are chosen to exercise different layers of the framework.  They
are not substitutes for the full checked Lean developments proposed in
\cref{sec:lean}; they show the complete human proof architecture and the exact
object endpoints.

\subsection{Least divisor as a named bounded certificate}

Let \(z\ge2\).  Define the bounded predicate
\[
  V(z,d)\;:\!\Longleftrightarrow\;
  2\le d\ \wedge\ d\le z\ \wedge\ \exists q\le z\;(d q=z).
\]
All quantifiers are bounded and all atoms are polynomial equalities or order
relations, so \cref{thm:reflection} supplies a Boolean verifier.  The value
\(d=z,q=1\) proves an accepted certificate below \(z+1\).  Apply
\cref{thm:least-code} to obtain the least accepted \(d^*(z)\).

The object theorem first proves that \(d^*\) divides \(z\) and is at least two.
To show primality, suppose \(d^*=ab\) with \(1<a<d^*\) and \(1<b\).  Then
\(a\) divides \(d^*\) and hence divides \(z\); the quotient can be composed
from the witnesses.  Thus \(a\) is an accepted divisor smaller than \(d^*\),
contradicting leastness.  Therefore \(d^*\) is prime.

The public term is justified because the certificate space has the explicit
bound \(z+1\).  If the theorem needed only ``some prime divides \(z\),'' the
certificate could instead be sealed after the primality proof.

\subsection{Rational bisection for the inverse of \texorpdfstring{\(y^3+y\)}{y cubed plus y}}

Fix an integer or rational input \(x\) and define
\(f_x(y)=y^3+y-x\).  The function is strictly increasing because
\[
 f_x(v)-f_x(u)
  =(v-u)(v^2+uv+u^2+1),
\]
and the second factor is positive when \(u<v\).  Choose rational endpoints
\(L_0,U_0\) with \(f_x(L_0)\le0\le f_x(U_0)\).  Such crude bounds are easy to
compute from \(|x|+1\).

At stage \(t\), let \(M_t=(L_t+U_t)/2\).  Exact rational sign comparison is a
bounded integer calculation after clearing positive denominators.  Set
\[
 (L_{t+1},U_{t+1})=
 \begin{cases}
   (M_t,U_t),& f_x(M_t)\le0,\\
   (L_t,M_t),& f_x(M_t)>0.
 \end{cases}
\]
Encode the pair of endpoints and common denominator in one state term.  The
reflected branch bit makes the step primitive recursive.

Use the quantifier-free invariant
\[
  f_x(L_t)\le0\ \wedge\ 0\le f_x(U_t)
   \ \wedge\
  2^t(U_t-L_t)=U_0-L_0.
\]
The base case is the initial bracket.  In the successor case, monotonicity and
the reflected sign branch preserve the bracket, while elementary rational
arithmetic proves the width equation.  \cref{thm:trace} gives the invariant at
an arbitrary query \(k\).  The midpoint or interval is therefore an
approximation to the unique inverse value, with an explicit error
\((U_0-L_0)2^{-k-1}\).

This example has no noncanonical certificate: the trace is deterministic.  It
exercises reflected branching, a packed state, query-first execution, and a
quantifier-free invariant.  A regular name follows by choosing a computable
number of bisection steps for each requested precision.

\subsection{Root clusters instead of unstable branch labels}

Suppose a finite algebra certificate produces disks
\(D_1,\ldots,D_m\) with multiplicities whose union contains all roots of a
polynomial and whose boundary tests certify the root counts.  At a collision,
individual labels are unstable, but the following observations are stable:

\begin{itemize}
\item the multiset of cluster centers and radii modulo permutation;
\item the multiplicity of each connected cluster;
\item symmetric polynomials of roots in a cluster;
\item the product factor associated with a cluster.
\end{itemize}

Let \(V(p,c)\) verify a cluster certificate and \(O(p,c)\) be one of these
observations.  A refinement theorem proves that any two valid certificates
have a common refinement and that their observations agree modulo permutation
or lie within a stated error.  \cref{thm:observed-relation} then exports a
coherent relational observation without choosing a branch.

If a separation residual
\[
  \operatorname{dist}(D_i,D_j)\ge\delta>0\quad(i\neq j)
\]
is supplied, request enough coefficient precision to verify the gaps and order
the disks lexicographically.  A bounded selector may then name each branch.
The label theorem is explicitly conditional on \(\delta\); it does not pretend
to pass continuously through a collision.

\subsection{Finite-net approximate minimization}

Let \(K\) be represented at precision \(r\) by a finite list
\(N_r=(x_0,\ldots,x_{M-1})\), a coverage relation saying every point of \(K\)
is within \(2^{-r}\) of some \(x_i\), and let \(f\) have a supplied modulus
\(\omega\).  Evaluate rational approximations \(v_i\) with error at most
\(2^{-s}\), and select the least index \(i^*\) minimizing the finite value
codes.

The finite search and comparison proof gives
\(v_{i^*}\le v_i\) for all net indices.  For any \(x\in K\), open the coverage
certificate to obtain a nearby \(x_j\).  The modulus gives
\(|f(x)-f(x_j)|\le\omega(2^{-r})\), and approximation errors compare
\(f(x_j)\) with \(v_j\) and \(v_{i^*}\) with \(f(x_{i^*})\).  Hence
\[
  f(x_{i^*})\le f(x)+2^{1-s}+\omega(2^{-r}).
\]
Seal the noncanonical nearby-net witness.  The public endpoint names the least
finite minimizer only because the list is explicitly bounded.  The coverage
and modulus formulas remain residuals unless separately produced.

\subsection{Conditional complete enumeration of Diophantine points}

Let \(P(\vec x)=0\) be an integer polynomial equation and let \(h(\vec x)\) be
a coded height.  Suppose an external mathematical argument supplies the object
producer
\[
  \Gamma\turnstilep
  \forall\vec x\,[P(\vec x)=0\to h(\vec x)\le H].
  \tag{HB}
\]
There is a primitive-recursive enumeration of every normalized integer or
rational tuple of height at most \(H\).  Filter it using reflected polynomial
equality and remove duplicates by reflected code equality.  Let \(L_H\) be the
resulting list.

A finite trace proves that each retained element is a solution and each
rejected enumerated code is either malformed or not a solution.  Given an
arbitrary solution \(x\), apply (HB) to place it below the enumeration bound;
normalization yields its unique code, and the filter proof shows that code is
retained.  Thus
\[
  \mathsf{HeightBound}(P,H),\Gamma
    \turnstilep \mathsf{CompleteList}(P,L_H).
\]

The finite enumeration theorem is unconditional; completeness is conditional
on the exact height producer.  Mere finiteness of the solution set cannot
replace (HB).

\subsection{A whole-path example}

Consider three finite stages.  At stage zero, choices \(a\in\{0,1\}\) are both
locally valid.  From \(a=0\), the only stage-one choice eventually dead-ends;
from \(a=1\), an accepted terminal path exists.  A greedy ``least valid
choice'' incorrectly selects \(0\).

Backward viability marks the terminal accepted nodes, then marks exactly the
stage-one nodes with accepted successors, and finally marks only \(a=1\) as
viable.  Forward least selection among \emph{viable} successors constructs the
accepted path.  The object proof is a fixed finite instance of the viability
theorem.  This tiny example captures the general reason separate existential
witnesses cannot simply be concatenated.

\section{Lean formalization architecture}
\label{sec:lean}

\subsection{Imports and namespace}

The reusable core should live beside \code{IPRADerivations}, not inside one
proof packet.  A minimal first import closure is

\begin{lstlisting}
import IPRADerivations.Core.Search
import IPRADerivations.Core.Pairing
import IPRADerivations.Core.Definitions

namespace LeanProofs.IPRA.SCOPE
open LeanProofs.ConciseIPRA LeanProofs.IPRA Syn
open scoped Fm
\end{lstlisting}

The generic library should not import Foundation, Busy Beaver, Diophantine
Consistency, or theorem-specific packets.  The open bridge is used only by a
separate export module after ordinary endpoints are stable.

\subsection{Structural kernel}

The accompanying \code{lean/SCOPE/Kernel.lean} records source-oriented wrappers
around checked rules:

\begin{lstlisting}
structure Contract (n : Nat) (Gamma : List (Fm PTm n))
    (goal : Fm PTm n) where
  valid : Fm PTm (n + 1)
  extract : valid :: Gamma.map Fm.shift
      |- [PRAAx] goal.shift

namespace Contract

theorem close (C : Contract n Gamma goal)
    (producer : Gamma |- [PRAAx] Fm.ex C.valid) :
    Gamma |- [PRAAx] goal :=
  .exE producer C.extract

theorem conditional (C : Contract n Gamma goal) :
    Gamma |- [PRAAx] (Fm.ex C.valid) --> goal := by
  apply Prf.impI
  apply Prf.exE Prf.hyp0
  exact C.extract.weaken (by simp)
end Contract
\end{lstlisting}

The file also wraps \code{Prf.trans}, \code{Prf.subst}, weakening, and boundary
discharge.  It contains no placeholder proof.  The first formalization
milestone is simply to elaborate it at the pinned revision and add it to the
axiom audit.

\subsection{Atomic schemas and bounded codes}

The schema-first interface is represented by:

\begin{lstlisting}
structure AtomSchema (k : Nat) where
  formula : Fm PTm k
  bit : PR k
  boolean : ([] : List (Fm PTm k)) |- [PRAAx]
    (app bit var == zero) \/ (app bit var == one)
  one_iff : ([] : List (Fm PTm k)) |- [PRAAx]
    Fm.iff (app bit var == one) formula

inductive BCode : Nat -> Type
  | atom  : AtomSchema k -> BCode k
  | conj  : BCode k -> BCode k -> BCode k
  | disj  : BCode k -> BCode k -> BCode k
  | impl  : BCode k -> BCode k -> BCode k
  | neg   : BCode k -> BCode k
  | exLt  : PTm k -> BCode (k + 1) -> BCode k
  | allLt : PTm k -> BCode (k + 1) -> BCode k
\end{lstlisting}

The actual archive uses repository notation in the definition of
\code{BCode.denote}.  An atom package stores only truth reflection; a generic
lemma derives falsity reflection from Booleanity.

\subsection{Arity-generic search}

The critical arithmetic implementation task is a parameterized version of
\code{Core/Search.lean}.  A suitable API is:

\begin{lstlisting}
namespace SearchN

def muF (body : PTm (k + 1)) : PR (k + 1)
def mu (body : PTm (k + 1)) (N : PTm n)
    (xs : Fin k -> PTm n) : PTm n

theorem search_spec (N i : PTm n) (xs : Fin k -> PTm n) :
  Gamma |- [PRAAx]
    le (mu body N xs) N /\
    ((lt (mu body N xs) N -->
       subst (Fin.cases (mu body N xs) xs) body == zero) /\
     (lt i (mu body N xs) -->
       neg (subst (Fin.cases i xs) body == zero)))

end SearchN
\end{lstlisting}

The exact vector expressions will use the repository's \code{Fin.cases} and
substitution conventions.  The proof should be obtained by parameterizing the
existing \code{spec\_open}; no new mathematical idea is required.  A useful
engineering strategy is to refactor the current two-argument search theorem as
the \(k=1\) corollary of \code{SearchN}.

\subsection{The reflection record}

The target is:

\begin{lstlisting}
structure Reflection {k : Nat} (A : BCode k) where
  char : PR k
  boolean : ([] : List (Fm PTm k)) |- [PRAAx]
    (app char var == zero) \/ (app char var == one)
  one_iff : ([] : List (Fm PTm k)) |- [PRAAx]
    Fm.iff (app char var == one) A.denote
  zero_iff : ([] : List (Fm PTm k)) |- [PRAAx]
    Fm.iff (app char var == zero) (neg A.denote)

def reflect (A : BCode k) : Reflection A
\end{lstlisting}

For atoms, return the stored symbol and proofs.  For connectives, build symbols
by \code{PR.comp} and proof terms by finite case analysis.  At bounded
quantifiers, construct a \code{SearchN.muF} symbol from the body characteristic
symbol and compose it with the bound's \code{PTm.toPR} symbol.

The substitution theorem should be almost definitional:

\begin{lstlisting}
@[simp] theorem char_subst (A : BCode k)
    (sigma : Fin k -> PTm n) :
  Syn.subst sigma (app (reflect A).char var) =
    app (reflect A).char sigma := by
  rfl
\end{lstlisting}

Depending on elaboration, the right proof may be \code{rfl} or the existing
\code{PTm.subst\_app} simplifier.  The main point is that no recursive theorem
about recompilation appears.

\subsection{Observer packages}

Exact observer erasure should be implemented first:

\begin{lstlisting}
structure ObserverContract (n : Nat) where
  valid : Fm PTm (n + 1)
  obs : PTm (n + 1)
  spec : Fm PTm (n + 1)
  extract : valid :: Gamma.map Fm.shift |- [PRAAx] spec
  coherent : valid.shift :: valid.reindexSecond ::
      Gamma.map (fun A => A.shift.shift) |- [PRAAx]
      obsFirst == obsSecond
\end{lstlisting}

The exact de Bruijn reindexing helpers should be defined once and tested on
small examples.  A public \code{out} formula existentially binds the
certificate and equates a public output variable with \code{obs}.  The module
then proves output existence and uniqueness.  Quotient observers replace
\code{==} by a reflected relation; metric observers add a rational distance and
triangle-inequality library.

\subsection{Indexed plan compiler}

Implement \code{Plan} only after the structural kernel and reflection module
compile.  Its first constructors are exact, specialize, supply, seal, and
boundary.  A simple structural function returns \code{Prf}.  The graph DSL is
strictly secondary: it should generate a \code{Plan}, not bypass it.

The effect checker is ordinary decidable host code.  It rejects:

\begin{itemize}
\item an unknown dependency or cycle not packaged as a trace node;
\item a missing producer with no named residual;
\item a relational output without a coherence derivation;
\item a named output without a direct producer or finite bound;
\item a sealed node whose certificate-dependent term is requested publicly.
\end{itemize}

Soundness of the checker is not trusted as an axiom.  Even if it accepts an
invalid plan, the invalid plan cannot produce a well-typed final \code{Prf}.

\subsection{Trace and demand modules}

A generic trace package stores fixed terms and local derivations:

\begin{lstlisting}
structure TraceSystem (k : Nat) where
  init : PTm k
  step : PTm (k + 2)
  inv : Fm PTm (k + 2)
  invQF : inv.QF
  base : ...
  next : ...
\end{lstlisting}

Use \code{recDef} for the state and \code{IPRA.ind} for the invariant theorem.
For capped demand closure, begin with fixed dimension at the metalevel.  Store
monotonicity as an object derivation, not as a semantic function property.  The
rank, inflation, stabilization, and leastness proofs are quantifier-free.

\subsection{Module layout}

\begin{verbatim}
IPRADerivations/Framework/SCOPE/
  Kernel.lean             -- structural wrappers
  Bits.lean               -- Boolean PR symbols and reflection
  SearchN.lean            -- arbitrary parameter arity
  BoundedCode.lean        -- schema, denotation, compiler
  Selection.lean          -- least accepted code
  Observer.lean           -- exact/quotient/metric output
  Trace.lean              -- deterministic traces
  Demand.lean             -- capped cyclic closure
  Viability.lean          -- whole finite paths
  Plan.lean               -- indexed proof plan and compiler
  Normalize.lean          -- host graph normalization
  Examples/
    LeastDivisor.lean
    CubicBisection.lean
    StableObserver.lean
    HeightList.lean
IPRABridge/SCOPE.lean      -- optional backend corollaries
\end{verbatim}

\subsection{Formalization order and tests}

The recommended order is:

\begin{enumerate}[label=\textbf{M\arabic*.}]
\item elaborate the structural kernel;
\item refactor arity-generic search;
\item prove bit combinators and atomic equality/order schemas;
\item implement bounded reflection and nested-quantifier tests;
\item add least selection;
\item add exact relational observers;
\item implement the indexed plan compiler;
\item add traces, demand closure, and viability;
\item migrate four pilot theorems;
\item expose bridge corollaries and run the repository axiom audit.
\end{enumerate}

Every final theorem module must contain no \code{sorry}, no theorem-specific
\code{axiom}, and a literal ordinary endpoint.  Tests should substitute
nontrivial compound terms under two alternating bounded binders; this is the
most direct regression test for the schema-first design.

\section{Repository integration and research program}
\label{sec:integration}

\subsection{Pilot migrations}

A useful migration sequence maximizes architectural coverage while minimizing
new mathematics.

\begin{enumerate}[label=(\arabic*)]
\item \textbf{Checked finite arithmetic.}  Wrap an existing finite-sum or
  bounded-search theorem as a sealed contract.  This tests contexts,
  specialization, and cut without adding arithmetic.
\item \textbf{Least divisor.}  Exercise bounded reflection and named selection.
\item \textbf{Cubic bisection.}  Exercise reflected branching, packed state,
  trace induction, and a regular-name endpoint.
\item \textbf{Stable observer.}  Use two different valid finite certificates
  with the same quotient observation; prove relational output uniqueness.
\item \textbf{Height-list theorem.}  Complete the verifier and enumeration while
  deliberately retaining one nontrivial height-bound residual.
\end{enumerate}

Only after these compile should broad packet migration begin.  The existing
reports remain useful provenance and mathematical design documents; the new
library theorem becomes the canonical formal endpoint.

\subsection{Acceptance criteria for a migrated packet}

A migrated packet should record:

\begin{itemize}
\item exact theorem names and source paths;
\item the ordinary \code{Prf PRAAx} endpoint;
\item the list of residual producer formulas;
\item the escape mode of each certificate;
\item whether each algorithmic component is recomputed or certified;
\item the explicit \code{QF} witness for every induction invariant;
\item the result of \texttt{\#print axioms};
\item a clear separation between kernel status, human proof status, and
  executable test status.
\end{itemize}

\subsection{Further research questions}

\begin{enumerate}[label=\textbf{Q\arabic*.}]
\item Can the arity-generic search proof be factored so that the current
  \code{Search} module becomes a one-line specialization, reducing maintenance
  rather than duplicating it?
\item What is the smallest useful library of reflected atomic schemas beyond
  equality and order: divisibility, bit extraction, rational comparison,
  membership in coded lists, and polynomial sign?
\item Can proof-plan normalization be verified by returning both a normalized
  plan and an equality/transport theorem, or is kernel checking of the final
  \code{Prf} sufficient for the desired trust boundary?
\item Which packet families can share one observational setoid library---root
  clusters, projective directions, quotient charts, or finite presentations?
\item Can metric observer composition be organized as an error-semiring so that
  precision demand propagation and triangle-inequality proofs are generated
  together?
\item For which metastable theorems can the repository prove additional
  summable-motion or monotonicity hypotheses that genuinely yield regular
  names?
\item Which qualitative finiteness theorems in the corpus admit explicit height
  producers in weak arithmetic, and which should remain conditional endpoints?
\item Can the native clause calculus substantially reduce generated proof size
  while preserving a small ordinary export theorem and transparent source
  provenance?
\item How much of the current packet corpus can be classified automatically by
  escape effect and residual type from its metadata?
\item Can one derive a general conservativity theorem saying that the SCOPE
  compiler introduces no principle beyond \(\PRAAx\), formalized entirely as
  the existence of its returned \code{Prf} value?
\end{enumerate}

\subsection{Conclusion}

The general obstacle to mathematics in iPRA is not that every concept must be
reduced to one enormous number.  The obstacle is interface discipline.  Finite
verification must be reflected into quantifier-free data; witnesses must be
scoped rather than silently chosen; outputs must respect nonuniqueness and
continuity; and deep global existence information must remain visible as a
producer boundary.

\Scope turns those principles into a concrete architecture.  Formula schemes
compile to fixed PR symbols, so open substitution is inherited from the
checked syntax.  Certificates are sealed, observed relationally, or named only
under justified selection.  A small indexed plan compiler assembles ordinary
object derivations, while trace, demand, and viability modules handle the
finite dynamics that recur throughout the repository.  The result is broad
enough for analysis, algebra, combinatorics, probability, topology, and
Diophantine applications, yet narrow enough that the next Lean work consists
of identifiable reusable modules rather than another informal universe.

\appendix
\section{Theorem and artifact ledger}
\label{app:ledger}

\small
\begin{longtable}{@{}p{0.30\textwidth}p{0.43\textwidth}p{0.19\textwidth}@{}}
\toprule
Result & Evidence & Status \\
\midrule
\endfirsthead
\toprule
Result & Evidence & Status \\
\midrule
\endhead
Ordinary \(\Prf\;\PRAAx\) calculus, substitution, weakening, cut & Existing
\code{ConciseIPRA/Deduction.lean} & Upstream checked \\
Quantifier-free induction and arbitrary-term instantiation & Existing
\code{IPRADerivations/Core/Basic.lean} & Upstream checked \\
Open-term closure to PR symbols and generic recursion & Existing
\code{Core/Definitions.lean} & Upstream checked \\
Bounded least search with one parameter & Existing
\code{Core/Search.lean} & Upstream checked \\
Scoped realization and conditional boundary & Direct use of \code{Prf.exE} and
\code{Prf.impI}; source candidate in archive & Human complete / candidate \\
Parameterized least-zero search & Full reduction to existing search proof with
arbitrary parameters & Human complete / Lean work \\
Schema-first bounded reflection & Full construction and proof in
\cref{sec:reflection,sec:connectives,sec:quantifiers} & Human complete / blueprint \\
Open substitution coherence & Fixed-symbol argument plus checked syntax laws &
Human complete \\
Least accepted certificate & Search and reflection proof in
\cref{thm:least-code} & Human complete / blueprint \\
Exact relational observer theorem & Nested existential elimination and
coherence & Human complete \\
Quotient/metric observer variants & Same scope proof with reflected relation or
distance & Human complete schema \\
Scope-minimal normal form & Finite commuting-conversion argument & Human complete \\
Indexed plan compiler & Structural recursion into checked proof constructors &
Human complete / blueprint \\
Deterministic trace invariant & Direct quantifier-free induction & Human complete \\
Capped cyclic demand closure & Rank and least-prefixed-point proof & Human complete;
reference tested \\
Backward viability and whole path & Bounded reflection and least search & Human
complete schema \\
Python reference algorithms & Sixteen standard-library tests & Executable only \\
Structural Lean kernel & No placeholders; source-audited against pinned API & Not
elaborated here \\
Bounded-code and compiler Lean files & Definitions and exact target signatures &
Blueprint, not checked \\
General height/modulus/cover producers & Theorem-specific, not supplied by the
framework & Explicit residuals \\
Exact selectors through degeneracy & Deliberately rejected without gap/producer &
Barrier; reformulated \\
\bottomrule
\end{longtable}
\normalsize

\section{Pinned source audit}
\label{app:sources}

The analysis and Lean signatures were pinned to ProveIt revision
\begin{center}
\code{8dfe26af18310f58aaff695fb0ce03ff5007d096}
\end{center}
on 16 November 2026.  The following files were used directly.

\begin{longtable}{@{}p{0.47\textwidth}p{0.46\textwidth}@{}}
\toprule
Repository path & Role in this report \\
\midrule
\endfirsthead
\toprule
Repository path & Role in this report \\
\midrule
\endhead
\code{Logic/PrimitiveRecursiveArithmetic/README.md} & canonical architecture and
status of packet versus checked layers \\
\code{Logic/PrimitiveRecursiveArithmetic/UNIFICATION-PLAN.md} & ordinary
open-sequent target, bridge plan, and migration priorities \\
\code{Logic/PrimitiveRecursiveArithmetic/ENDPOINTS.md} & pinned build endpoints
and bridge status \\
\code{ConciseIPRA/Lean/ConciseIPRA/Deduction.lean} & exact constructors and
structural theorem signatures \\
\code{ConciseIPRA/Lean/ConciseIPRA/PR.lean} & \(\PR\), \(\PTm\),
\(\PRAAx\), and induction \\
\code{Derivations/Lean/IPRADerivations/Core/Basic.lean} & \code{inst},
\code{ind}, and arithmetic terms \\
\code{Core/Definitions.lean} & \code{PTm.toPR}, \code{recDef}, sums and products \\
\code{Core/Order.lean} & quantifier-free order and decidability results \\
\code{Core/Pairing.lean} & equality bit, pairing, and simultaneous recursion \\
\code{Core/Search.lean} & least-zero implementation and proof pattern \\
\code{docs/proofs/README.md} & survey of proof packets and their status \\
\code{docs/proofs/foundations/FINER-iPRA-framework-2026-11-16/} & finite-query
compiler, demand, viability, observer, and difficulty analysis \\
\code{docs/proofs/foundations/scoped_certificate_realization_ipra_2026-11-16/}
& certificate escape analysis and structural kernel candidate \\
\code{docs/proofs/foundations/ipra_bounded_reflection_2026-11-16/} & earlier
bounded-reflection plan and its open obligations \\
\bottomrule
\end{longtable}

The present report does not claim that the accompanying Lean blueprint was
compiled.  The runtime used for artifact generation had \LaTeX{} and Python but
no \code{lean} or \code{lake} executable.  Repository integration must repeat
the source audit at the then-current revision and run focused builds plus
\texttt{\#print axioms}.

\section{Executable reference artifacts}
\label{app:reference}

The archive's \code{reference/scope\_reference.py} is a standard-library model
of the finite algorithms.  It includes:

\begin{itemize}
\item an AST and evaluator for equality, order, Boolean connectives, and nested
  bounded quantifiers;
\item least bounded witness search and a minimality certificate;
\item monotone capped demand closure with a runtime rank check;
\item a finite DAG checker for sealed, relational, and named escape modes;
\item exact observer independence and metric diameter checks;
\item a JSON demonstration generator.
\end{itemize}

The bundled unit suite has sixteen tests.  It includes two nested bounded
quantifiers, implication truth tables, parameter environments under binders,
least search with and without a witness, a cyclic demand system, rejection of
an unjustified exact name, a missing observer-coherence proof, and exact and
metric observer examples.

Run from the archive root:

\begin{verbatim}
python3 -m unittest -v reference/test_scope_reference.py
python3 reference/scope_reference.py --pretty > examples/demo.json
\end{verbatim}

Passing these tests shows that the reference algorithms and metadata checks
behave as documented on the examples.  It does not prove the universal iPRA
reflection or compiler theorems.  Those are certified only when the Lean
modules construct accepted \code{Prf PRAAx} terms.


\end{document}
